\section{Introspection Games}
\label{sec:introspection}

\newcommand{\ai}[1][y,\, a]{A^\Intro_{#1}}
\newcommand{\bs}[1][z,\, a]{B^\Sample_{#1}}
\newcommand{\br}[1][y,y^\perp, a]{B^\Read_{#1}}
\newcommand{\ahk}[1][y_{<k}]{A^{\Hide{k}}_{#1}}
\newcommand{\ahp}[1][y_{\le k}]{A^{\Hide{k+1}}_{#1}}
\newcommand{\bhk}[1][y_{<k}]{B^{\Hide{k}}_{#1}}
\newcommand{\bhp}[1][y_{\le k}]{B^{\Hide{k+1}}_{#1}}
\newcommand{\z}[1][z]{\sigma^Z_{#1}}
\newcommand{\x}[1][x]{\sigma^X_{#1}}

\newcommand{\La}{[L(\cdot) = y]}
\newcommand{\Lk}[1][k]{[L_{#1,\, y_{<#1}}(\cdot) = y_{#1}]}
\newcommand{\Llk}[1][k]{[L_{<#1}(\cdot) = y_{<#1}]}
\newcommand{\Llek}[1][k]{[L_{\le #1}(\cdot) = y_{\le #1}]}
\newcommand{\Lperpk}[1][k]{[L_{#1,\, y_{<#1}}^\perp (\cdot) = y_{#1}^\perp]}

\newcommand{\aigek}[1][y_{\ge k},\, a]{A^{\Intro,\, y_{<k}}_{#1}}
\newcommand{\aigk}[1][y_{> k},\, a]{A^{\Intro,\, y_{\le k}}_{#1}}
\newcommand{\bsL}{\bs[\La,\, a]}
\newcommand{\zL}{\z[\La]}
\newcommand{\zLk}[1][k]{\sigma^{Z}_{\Lk[#1]}}
\newcommand{\zLlk}[1][k]{\z[{\Llk[#1]}]}
\newcommand{\zLlek}{\z[\Llek]}
\newcommand{\xLperpk}[1][k]{\sigma^{X}_{\Lperpk[#1]}}

\newcommand{\aizk}[1][y,\, a]{A^{\Intro,\, Z_{<k}}_{#1}}

\subsection{Overview}
\label{sec:intro-overview}

Consider a normal form verifier $\verifier = (\sampler, \decider)$ (see
Definition~\ref{def:normal-game}).
In this section we design a normal form verifier $\verifier^\intro =
(\sampler^\intro, \decider^\intro)$ (called the \emph{introspective verifier})
such that in the $n$-th game $\verifier^\intro_n$ (called the
\emph{introspection game}; see Definition~\ref{def:normal-game} for the
definition of the $n$-th game associated with a normal form verifier) the
verifier expects the players to sample for themselves questions $x$ and $y$
distributed as their questions in the game $\verifier_N$ for index $N =
2^n$---this is the ``introspection'' step.
Note the exponential separation of the indices of the introspection game versus
the original game!
Our construction of the introspection game generalizes the introspection
technique of~\cite{NW19}.

Recall the execution of the $N$-th game corresponding to the ``original''
verifier $\verifier$ (see Definition~\ref{def:normal-game}).
Let $n\geq 1$, $N=2^n$, and suppose that the CL functions of $\sampler$ on index
$N$ are $L^\alice,L^\bob$ acting on an ambient space $V$.
In the game $\verifier_N$, the verifier first samples $z \in V$ uniformly at
random and gives each player $w\in\AB$ the question $x_w = L^w(z)$.
In this context, the string $z$ is referred to as the ``seed''.
The players respond with answers $a_\alice$ and $a_\bob$, respectively, and the
verifier accepts or rejects according to the output of $\decider(N, x_\alice,
x_\bob, a_\alice, a_\bob)$.

In the introspection game, with some constant probability independent of $n$,
the verifier $\verifier^\intro_n$ sends the question pair $(\Introspect,\alice)$
to player $w$ and $(\Introspect,\bob)$ to the other player $\overline{w}$, where
$w \in \AB$ and recall that $\overline{w}=\bob$ if $w=\alice$ and
$\overline{w}=\alice$ otherwise.
The verifier expects player $w$ to measure their share of the state
$\ket{\EPR}_V$ using a coarse-grained $Z$-basis measurement whose outcomes range
over $L^\alice(V)$, and similarly player $\overline{w}$ measures the state
$\ket{\EPR}_V$ using a coarse-grained $Z$-basis measurement with outcomes that
range over $L^\bob(V)$.
If the players perform these measurements honestly, then the outcomes
$(x_\alice,x_\bob)$ are distributed exactly according to $\mu_{\sampler,\, N}$, the
question distribution of the game $\verifier_N$.
Players $w$ and $\overline{w}$ are then expected to respond with the question
$x_w$ and $x_{\ol{w}}$ that they each sampled, together with strings $a_w$ and
$a_{\overline{w}}$, respectively, which are intended to be the answers for the
question pair $(x_\alice,x_\bob)$ in the game $\verifier_N$.
In other words, the players introspectively sample the question pair
$(x_\alice,x_\bob)$ and then respond with the question itself and an answer for
it.

To facilitate comprehension, we call the players that interact with the
introspective verifier $\verifier^\intro$ the ``introspecting players'', and the
players that interact with the ``original'' verifier $\verifier$ the ``original
players''.
To ensure that the introspecting players follow the above procedure honestly,
the introspective verifier $\verifier^\intro_n$ first uses the (binary) Pauli
basis test described in Section~\ref{sec:pauli-verifier} to force the
introspecting players to share the state $\ket{\EPR}_{V}$.
The Pauli basis test also ensures that the players measure $\sigma^W$ and report
the outcome honestly when they receive questions $(\Pauli, W)$ for $W \in \{X,
Z\}$.
For $v\in\{\alice,\bob\}$ the verifier cross-checks the question pairs
$(\Introspect, \trole)$ and $(\Pauli, Z)$ to enforce that the honest measurement
is performed for question $(\Introspect, \trole)$.

The main difficulty in the soundness analysis is to ensure that the answer of
player $w$ who received question $(\Introspect, \trole)$ depends only on
$L^\trole(z)$ and not on any other information about the string $z$.
First assume for simplicity that $L^\trole$ is a linear function.
As shown below (based on Lemma~\ref{lem:why-didnt-i-think-of-this-before}),
$L^\trole(z)$ can be obtained by measuring a specific collection of~$\sigma^Z$
observables; namely, the set
\begin{equation}
  \label{eq:intro-overview-z}
  \{\sigma^Z(u) : u \in \ker(L^\trole)^\perp\}.
\end{equation}
To prevent the player from obtaining any additional information the verifier
needs to enforce that the player does \emph{not} additionally measure any
$\sigma^Z(u)$ for $u \notin \ker(L^\trole)^\perp$.
(We refer to any such $\sigma^Z$ as a ``prohibited'' $\sigma^Z$.)

The introspective verifier achieves this by sometimes sending the $\Read$ type
question $(\Read, \trole)$ or $\Hide{k}$ type questions $(\Hide{k}, \trole)$ to the
players.
When receiving the $\Read$ question, the players are required to also measure
observables from the set
\begin{equation}
  \label{eq:intro-overview-x}
  \{\sigma^X(r) : r \in \ker(L^\trole)\},
\end{equation}
which (as shown in Lemma~\ref{lem:commute} below) commute with every $Z$-basis
measurement in~\eqref{eq:intro-overview-z}.
On the other hand, any prohibited $\sigma^Z(u)$ must anticommute with at least
one of the $\sigma^X(r)$, as otherwise $u$ would be in $\ker(L^\trole)^\perp$.
As a result, honestly measuring the~$\sigma^X$ observables
of~\eqref{eq:intro-overview-x} has the effect of preventing the player from
measuring any of the prohibited~$\sigma^Z$ observables, so that the answer $a$
can effectively only depend on the question $L^\trole(z)$.
(In the protocol, the verifier asks the player to measure the function
$(L^\trole)^\perp(z)$ in the $X$ basis, rather than all of the~$\sigma^X$
observables.
By Lemma~\ref{lem:why-didnt-i-think-of-this-before} the two are equivalent.)
Similarly to how questions $(\Pauli, Z)$ and $(\Introspect, \trole)$ are
cross-checked, the questions $(\Pauli, X)$, $(\Hide{k}, \trole)$ and $(\Read,
\trole)$ are cross-checked in order to ensure that the honest $X$-basis
measurements are performed for the $\Hide{k}$ questions and the $\Read$ questions.

When the CL functions $L^\trole$ are $\ell$-level for $\ell > 1$, the
introspective verifier sends one of $O(\ell)$ different hiding questions to the
players, chosen at random; together these hiding checks ensure that each of the
constituent linear maps of $L^\trole$ are honestly measured.
Intuitively, these hiding questions ``interpolate'' between questions
$(\Pauli,Z)$, $(\Introspect,\trole)$ and $(\Pauli,X)$ in a way that, for all
pairs of questions asked by the verifier, the honest measurements commute.
(See Figure~\ref{fig:intro-honest} for an overview of the honest measurements.)
This strong commuting property of the strategy is essential for the
oracularization transformation in \cref{sec:oracle} to be possible.

A key property of the introspection game is that the distribution of questions
(which include the Pauli basis test questions as well as the introspection
questions and the hiding questions) is also conditionally linear.
This means that the introspection game can be ultimately specified by a normal
form verifier $\verifier^\intro = (\sampler^\intro,\decider^\intro)$, which is
crucial for recursive compression of games.
Moreover, while the time complexity of the introspective verifier's decider
$\decider^\intro$ remains polynomially related to that of $\decider$ on index
$N$, the time complexity of the sampler $\sampler^\intro$ is $\polylog(N)$
(exponentially smaller), due to the efficiency of the Pauli basis test.
Finally, $\sampler^\intro$ only depends on $\verifier$ through the number of
levels $\ell$ of $\sampler$ as well as upper bounds on the time complexities of
$\sampler$ and $\decider$.

\subsection{The introspective verifier}
\label{sec:intro-verifier}

\def\tsint{\hat{\sampler}^\intro}
\def\tdint{\hat{\decider}^\intro}
\def\tvint{\hat{\verifier}^\intro}
\def\sint{\sampler^\intro}
\def\dint{\decider^\intro}
\def\vint{\verifier^\intro}

Let $\verifier = (\sampler, \decider)$ be an $\ell$-level normal form verifier.
We call $\verifier$, $\sampler$, and $\decider$ the ``original'' verifier,
sampler, and decider, respectively.
For all original verifier $\verifier$, we define the introspective verifier
$\vint$ in this section.
We use $N = 2^n$ to denote the index of the original verifier $\verifier$ that
is simulated by the introspective verifier $\vint$ on index $n$.





The \emph{introspective verifier corresponding to $\verifier$ and parameters
  $(\lambda,\ell)$} is a \emph{typed} normal form verifier $\tvint
=(\tsint, \tdint)$, sketched in
Section~\ref{sec:intro-overview} and specified in detail in the present section
(see Section~\ref{sec:types} for the definition of typed normal form verifiers).
In the following descriptions of the sampler $\tsint$ and decider
$\tdint$, all parameters are functions of the index $n$, the number of
levels $\ell$ of the sampler $\sampler$, and the parameter $\lambda$; we often
leave this dependence implicit.

Let $R = N^\lambda$, and let $\introparams(R) =
(q,m,d)$ denote the parameter tuple specified in
Section~\ref{sec:qld-complexity}.
Note that $\introparams$ is implicitly a function of $n$ (since $R$ is a
function of $n$).
The associated parameter tuple $\introparams$ is intended to parametrize a Pauli
basis test that certifies the \hnote{edited in response to LB comments:} presence of $Q = 2^m \log q$-qubit EPR states; intuitively, these qubits are meant to serve as the source of randomness for the CL functions of the
original sampler $\sampler$. Since we will assume that the original verifier $\verifier$ is $\lambda$-bounded, the original sampler $\sampler$ on index $N$ has time complexity at most $R$, and therefore the amount of randomness needed is at most $R$. By \Cref{lem:delta-bound}, the number of qubits $Q$ in the EPR states is at least $R$, and thus the EPR qubits can be used to properly introspect the original game. 

Recall that the players in the introspection game are referred to as
``introspecting players'' and the players in the original game are referred to
as ``original players''.
We use the following notation in order to distinguish between questions and
answers meant for the introspecting players versus the original players.
The questions and answers of the introspecting players are denoted by hatted
variables (e.g., $\hat{x}$ and $\hat{a}$).
Similarly, the associated question types are denoted using hatted variables
$\hat{\tvar}$.
The questions and answers of the original players in the original game
$\verifier_N$ are denoted using non-hatted variables (e.g.\ $x$, $a$, and so
on).

\paragraph{Types and type graph.}
The type set $\type^\intro$ for the introspective verifier $\tvint$ is
\begin{gather*}
  \type^\intro = \type^\pauli \cup \Big( \Big(\big\{ \Introspect, \Sample,
  \Read \big\} \cup \Big(\, \bigcup_{k=1}^{\ell} \bigl\{ \Hide{k} \bigr\} \,
  \Big) \Big) \times \ab \Big) \;,
\end{gather*}
where $\type^\pauli$ is the type set of the Pauli basis test, defined in
Section~\ref{sec:pauli-verifier}.
The type graph $G^\intro$ is specified in Figure~\ref{fig:type-graph-intro}.

\begin{figure}[!htbp]
  \centering
  \begin{tikzpicture}[scale=.8]

    \tikzset{type/.style args={[#1]#2}{
        draw,circle,fill,scale=0.25,
        label={[font=\scriptsize, label distance=1pt]#1:#2}
      }}



    \foreach \i in {1,...,6} \draw (0,8-9/8*\i)
    coordinate (Constraint-\i)
    node[type={[180]$\Constraint_\i$}] {};

    \foreach \i in {1,...,9} \draw (2.5,9-\i)
    coordinate (Variable-\i)
    node[type={[330]$\Variable_\i$}] {};

    \foreach \i in {1,...,3} \foreach \j in {1,...,3}
    \pgfmathsetmacro{\k}{(\i-1)*3+\j}
    \draw (Constraint-\i) -- (Variable-\k);

    \foreach \i in {4,...,6} \foreach \j in {1,...,3}
    \pgfmathsetmacro{\k}{\i-3+(\j-1)*3}
    \draw (Constraint-\i) -- (Variable-\k);



    \draw (6,9) coordinate (DLine-X) node[type={[90]$(\DLine, X)$}] {};
    \draw (8,9) coordinate (Plane-X) node[type={[90]$(\ALine, X)$}] {};
    \draw (7,8) coordinate (Point-X) node[type={[315]$(\Point, X)$}] {};
    \draw (9.5,8) coordinate (Pauli-X) node[type={[0]$(\Pauli, X)$}] {};
    \draw (6,3) coordinate (DLine-Z) node[type={[270]$(\DLine, Z)$}] {};    
    \draw (8,3) coordinate (Plane-Z) node[type={[270]$(\ALine, Z)$}] {};
    \draw (7,4) coordinate (Point-Z) node[type={[45]$(\Point, Z)$}] {};
    \draw (9.5,4) coordinate (Pauli-Z) node[type={[0]$(\Pauli, Z)$}] {};
    \draw (8,6) coordinate (Pair) node[type={[0]$\Pair$}] {};
    \draw (7.3,6.8) coordinate (Pair-X) node[type={[0]$(\Pair,X)$}] {};
    \draw (7.3,5.2) coordinate (Pair-Z) node[type={[0]$(\Pair,Z)$}] {};



    \draw (10.5,8.7) coordinate (Hide-1-A) node[type={[90]$(\Hide{1}, \alice)$}] {};
    \draw (12.5,8.7) coordinate (Hide-2-A) node[type={[90]$(\Hide{2}, \alice)$}] {};
    \draw (13.8,8.7) node {$\cdots$};
    \draw (15,8.7) coordinate (Hide-l-A) node[type={[90]$(\Hide{\ell}, \alice)$}] {};

    \draw (10.5,7.3) coordinate (Hide-1-B) node[type={[270]$(\Hide{1}, \bob)$}] {};
    \draw (12.5,7.3) coordinate (Hide-2-B) node[type={[270]$(\Hide{2}, \bob)$}] {};
    \draw (13.8,7.3) node {$\cdots$};
    \draw (15,7.3) coordinate (Hide-l-B) node[type={[270]$(\Hide{\ell}, \bob)$}] {};

    \draw (10.5,3.3) coordinate (Sample-A) node[type={[270]$(\Sample, \alice)$}] {};
    \draw (12.75,3.3) coordinate (Intro-A) node[type={[270]$(\Introspect, \alice)$}] {};
    \draw (15,3.3) coordinate (Read-A) node[type={[270]$(\Read, \alice)$}] {};

    \draw (10.5,4.7) coordinate (Sample-B) node[type={[90]$(\Sample, \bob)$}] {};
    \draw (12.75,4.7) coordinate (Intro-B) node[type={[90]$(\Introspect, \bob)$}] {};
    \draw (15,4.7) coordinate (Read-B) node[type={[90]$(\Read, \bob)$}] {};



    \foreach \from/\to in {DLine-X/Point-X, Plane-X/Point-X, Point-X/Pauli-X,
      DLine-Z/Point-Z, Plane-Z/Point-Z, Point-Z/Pauli-Z,
      Point-X/Variable-1, Point-Z/Variable-5,
      Hide-2-A/Hide-1-A,
      Hide-2-B/Hide-1-B,
      Sample-A/Intro-A, Intro-A/Read-A,
      Sample-B/Intro-B, Intro-B/Read-B,
      Intro-A/Intro-B}
    \draw (\from) -- (\to);

    \draw (Pauli-X) to [out=60,in=180] (Hide-1-A);
    \draw (Pauli-X) to [out=300,in=180] (Hide-1-B);
    \draw (Pauli-Z) to [out=60,in=180] (Sample-B);
    \draw (Pauli-Z) to [out=300,in=180] (Sample-A);
    \draw [rounded corners=8] (Hide-l-A) to +(2.2,0) to +(2.2,-5.4) to (Read-A);
    \draw [rounded corners=8] (Hide-l-B) to +(1.2,0) to +(1.2,-2.6) to (Read-B);
    \draw (Point-X) to [out=275,in=120] (Pair-X) to [out=300,in=140] (Pair);
    \draw (Point-Z) to [out=85,in=240] (Pair-Z) to [out=60,in=220] (Pair);

  \end{tikzpicture}
  \caption{Type graph $G^\intro$ for the introspection game.
    Each vertex also has a self-loop which is not drawn in the figure for
    clarity. }
  \label{fig:type-graph-intro}
\end{figure}

\paragraph{Sampler.}
We first define a $3$-level $(\type^\intro,G^\intro)$-typed sampler
$\tilde{\sampler}^\intro$, which has field size $q(n)$ and dimension $3m(n) +
3$, where $q(n)$ and $m(n)$ are specified by $\introparams(R)$.
\hnote{clarified sentence in response to Bowen comments} Note that the dimension matches that of the space $V^\pauli$ of the CL
functions of the Pauli basis test parameterized by $\introparams(R)$, specified in
Section~\ref{sec:qld-game}. 

Fix $n \in \N$.
We specify the CL functions of $\tilde{\sampler}^\intro$ on index $n$.
Since the functions $L_{\hat{\tvar}}^{\alice,\, n}$ and
$L_{\hat{\tvar}}^{\bob,\, n}$ are identical for all $n$ and $\hat{\tvar}$, we
omit the superscripts $\alice$ and $\bob$.
For types $\hat{\tvar} \in \type^\pauli$, the CL functions $L_{\hat{\tvar}}^n$
are given by those specified in Section~\ref{sec:qld-game}, parameterized by
$\introparams(R)$.
For types $\hat{\tvar} \in \type^\intro \setminus \type^\pauli$, the associated
CL functions are defined to be $0$-level CL functions (i.e., they are the $0$
map).
This means that for question types such as $\hat{\tvar}=(\Introspect, \trole)$
or $\hat{\tvar}=(\Sample, \trole)$ for $\trole\in\{\alice,\bob\}$ the associated question is solely comprised
of the question type label.

Finally, we define the typed sampler $\tsint$ to be
$\downsize(\tilde{\sampler}^\intro)$, the downsized sampler
(\cref{def:downsize-typed-sampler}) corresponding to $\tilde{\sampler}^\intro$.
The typed CL functions associated with
the Pauli basis test are $3$-level; using \cref{rk:higher-level,lem:cl-downsize}
it follows that $\tsint$ is a $3$-level typed sampler.

The following lemma establishes the complexity of the sampler $\tsint$
as well as the complexity of computing a {description} of it from the parameters
$(\lambda,\ell)$.
\begin{lemma}
  \label{lem:intro-sampler-complexity}
  There exists a $2$-input Turing machine $\ComputeIntroSampler$ that
  on input $(\lambda,\ell)$ outputs a description of the sampler
  $\tsint$ (corresponding to parameters $\lambda,\ell$) in time $\polylog(\lambda,\ell)$.
  Furthermore, 
  \begin{enumerate}
  \item $\TIME_{\tsint}(n) = \poly(n, \lambda, \ell)$, and
	\item $\tsint$  is a $3$-level typed sampler.
  \end{enumerate}
\end{lemma}

\begin{proof}
  Define the following $9$-input Turing machine
  $\RawIntroSampler$, that does not depend on any parameters (so that its description length is constant). 
  On input $(\lambda', \ell', n', x_1', \ldots, x_6')$, the Turing machine $\RawIntroSampler$
  computes the output of the typed sampler $\tsint$ (parameterized by
  $(\lambda, \ell)$) with input tapes set to $(n', x_1', \ldots, x_6')$.
  In more detail, the Turing machine $\RawIntroSampler$ first computes
  $\introparams(R)$ for $R= 2^{\lambda' n'}$.
  Using \cref{lem:introparams-complexity}, this computation takes time
  $\poly\log(R)$.
  Next, depending on the contents of the last $7$ input tapes of
  $\RawIntroSampler$, the Turing machine evaluates the dimension of
  $\tsint$ (which can easily be computed from $\introparams(R)$), or
  one of the CL functions, or returns a representation of a factor space of
  $\tsint$.
  If the type passed as input is $\hat{\tvar} \in \type^\pauli$ then by
  Lemma~\ref{lem:qld-complexity} this takes time $\polylog(R)$.
  If $\hat{\tvar} \in \type^\intro\backslash\type^\pauli$ then this can be done
  in $O(\log \ell')$ time (to read the input type).
  Overall, $\RawIntroSampler$ runs in time $\poly\log(R,\ell')$.
  
  We now define the Turing machine $\ComputeIntroSampler$ as follows.
  On input $(\lambda,\ell)$, $\ComputeIntroSampler$ outputs the description of
  $\RawIntroSampler$ with the first two input tapes hardwired to $\lambda$ and
  $\ell$, respectively, yielding the sampler $\tsint$ corresponding to
  parameters $(\lambda,\ell)$.
  Computing this description takes $O(\log \lambda + \log \ell)$ time.
  
  The time complexity of $\tsint$ follows from the time complexity of
  $\RawIntroSampler$ and the number of levels is by construction.
\end{proof}

\paragraph{Decider.}
The typed decider $\tdint$ is specified in
Figure~\ref{fig:intro-decider}.
We explain how to interpret the figure, including the notation.
(It may also be helpful to review the description of the intended strategy for
the players in the game, described in Section~\ref{sec:intro-completeness}.)

\begin{figure}[!htbp]
  \begin{gamespec}
    \begin{table}[H]
      \centering
      \small
      \begin{tabularx}{.7\textwidth}{l X}
        \toprule
        Type \hspace{6em} & Answer format\\
        \midrule
        $(\Introspect, \trole)$ & $(y,a) \in V \times \{0,1\}^*$ \\
        $(\Sample, \trole)$ &  $(z,a) \in V \times \{0,1\}^*$ \\
        $(\Read, \trole)$ & $(y,y^{\perp},a) \in V \times V \times \{0,1\}^*$\\
        $(\Hide{k}, \trole)$ & $(y,y^{\perp},x) \in V \times V \times V$\\
        \bottomrule
      \end{tabularx}
    \end{table}
    In the following, whenever $\sampler$ or $\decider$ is called as a subroutine,
    $\tdint$ aborts and rejects if the subroutine takes more than
    $N^\lambda$ time steps.
    On input $(n, \hat{\tvar}_\alice, \hat{x}_\alice, \hat{\tvar}_\bob,
    \hat{x}_\bob, \hat{a}_\alice, \hat{a}_\bob)$, the decider $\tdint$
    first computes the dimension $s(N)$ of $V$ by calling the original sampler
    $\sampler$ on input $(N,\gamestyle{dimension})$.
    The decider rejects if $s(N) > N^\lambda$, or if \hnote{edited:} $\max \bigl\{ \abs{\hat{a}_\alice},
    \abs{\hat{a}_\bob} \bigr\} \geq 3 \cdot 2^m \cdot \log q$ where $\introparams(N^\lambda) = (q,m,d)$. Otherwise, it performs the following tests for all $w,\trole \in \AB$.
    (If no test applies, the decider accepts.) 
    

\begin{enumerate}[itemsep=2pt, parsep=2pt]
    \item (\textbf{Pauli test}).
      \label{enu:pauli}
      If $\hat{\tvar}_\alice, \hat{\tvar}_\bob \in \type^\pauli$, the decider
      accepts if $\decider^\pauli_\introparams$ accepts the input
      $(\hat{\tvar}_\alice, \hat{x}_\alice, \hat{\tvar}_\bob, \hat{x}_\bob,
      \hat{a}_\alice, \hat{a}_\bob)$.
    \item (\textbf{Sampling tests}).
      \label{enu:sampling}
      \begin{enumerate}
      \item If $\hat{\tvar}_w = (\Pauli,Z)$ and $\hat{\tvar}_{\overline{w}} =
        (\Sample, \trole)$, accept if $\hat{a}_w^V = z_{\overline{w}}$.
        \label{enu:sampling-pauli}

      \item If $\hat{\tvar}_w = (\Introspect, \trole)$,
        $\hat{\tvar}_{\overline{w}} = (\Sample, \trole)$, accept if $y_{w} =
        L^\trole(z_{\overline{w}})$ and $a_w=a_{\overline{w}}$.
        \label{enu:sampling-intro}
      \end{enumerate}

    \item (\textbf{Hiding tests}).
      \label{enu:hiding}
      \begin{enumerate}

      \item If $\hat{\tvar}_w = (\Introspect, \trole)$,
        $\hat{\tvar}_{\overline{w}} = (\Read, \trole)$, accept if $y_{w} =
        y_{\overline{w}}$ and $a_w = a_{\overline{w}}$.
        \label{enu:hiding-intro}

      \item If $\hat{\tvar}_w = (\Hide{\ell}, \trole)$,
        $\hat{\tvar}_{\overline{w}} = (\Read, \trole)$, accept if $y_{w,\,
          <\ell} = y_{\overline{w},\, <\ell}$, and $y^{\perp}_w =
        y^{\perp}_{\overline{w}}$.
        \label{enu:hiding-read}
        
      \item \label{enu:hiding-same} If $\hat{\tvar}_w = (\Hide{k}, \trole)$,
        $\hat{\tvar}_{\overline{w}} = (\Hide{k+1}, \trole)$ for some $k \in \{ 1,
        2, \ldots, \ell-1 \}$, accept if
        \begin{equation*}
          y_{w,\, <k} = y_{\overline{w},\, <k}\;, \quad
          y^{\perp}_{w,\, \le k} = y^{\perp}_{\overline{w},\, \le k}\;, \quad
          x_{w,\, >k+1} = x_{\overline{w},\, >k+1}\;,
        \end{equation*}
        and $y^{\perp}_{\overline{w},\, k+1} = \big( L^{\trole}_{k+1,\, u} \big)^\perp
        (x_{w,\, k+1})$ where $u = y_{\overline{w},\, \le k}$.
      \item If $\hat{\tvar}_w = (\Pauli,X)$, $\hat{\tvar}_{\overline{w}} =
        (\Hide{1}, \trole)$, accept if $y^\perp_{\overline{w},\, 1} =
        \bigl( L^{\trole}_1 \bigr)^\perp (\hat{a}_w^{V_1})$ and $\hat{a}_w^{V_{> 1}} =
        x_{\overline{w},\, >1}$.
        \label{enu:hiding-pauli}
      \end{enumerate}

    \item (\textbf{Game test}).
      \label{enu:intro-game}
      If $\hat{\tvar}_w = (\Introspect, \alice)$ and $\hat{\tvar}_{\overline{w}}
      = (\Introspect, \bob)$, accept if $\decider$ accepts $(N, y_w,
      y_{\overline{w}}, a_w, a_{\overline{w}})$.

    \item (\textbf{Consistency test}).
      \label{enu:intro-consistency}
      If $\hat{\tvar}_\alice = \hat{\tvar}_\bob$, accept if and only if
      $\hat{a}_\alice = \hat{a}_\bob$.
    \end{enumerate}
  \end{gamespec}
  \caption{The typed decider $\tdint$ for the introspective verifier,
    parameterized by integers $\lambda, \ell$, on index $n$.
    $N$ denotes $2^n$, $V$ is the ambient space for $\sampler$, and
    $\{L^\trole\}_{\trole\in\{\alice,\bob\}}$ the associated CL functions on
    index $N$.}
  \label{fig:intro-decider}
\end{figure}

\medskip
\emph{Question and answer format.}
The decider takes as input a tuple $(n, \hat{\tvar}_\alice, \hat{x}_\alice,
\hat{\tvar}_\bob, \hat{x}_\bob, \hat{a}_\alice, \hat{a}_\bob)$ where
$(\hat{\tvar}_w,\hat{x}_w)$ denotes the question for introspecting player $w \in
\AB$, and $\hat{a}_w$ denotes their answer.
As in the specification of the Pauli basis test, in
Figure~\ref{fig:intro-decider} we include an ``answer key'' indicating how the
players' answers are parsed, depending on the question type.
When the question type is from $\type^\pauli$, the question and answer format
are as described in Figure~\ref{fig:decider_pauli}.
When the question type is in $\type^\intro \setminus \type^\pauli$, the answer
format is described in the table at the top of
Figure~\ref{fig:intro-decider}.\footnote{There is no question format
  specification for the question types in $\type^\intro \setminus \type^\pauli$,
  because the question is solely comprised of the question type label.}
For each such question, there is an associated variable $\trole \in \AB$ that
indicates to the introspecting player which {original player} it is supposed to
impersonate in the introspection game.

In the figure, $V$ denotes the ambient space of the original sampler $\sampler$
on index $N = 2^n$. \hnote{added:} We assume that the original verifier $\verifier$ is $\lambda$-bounded and in particular the running time of the original sampler $\sampler$ (and therefore its dimension) is at most $R = N^\lambda$. If the dimension were larger than $N^\lambda$, the decider $\hat{\decider}^\intro$ would always reject in the beginning. 

Since $V$ is isomorphic to $\F_2^{s(N)}$, where $s(N)$ is the dimension of
$\sampler$, the space $V$ is identified in a canonical way as the register
subspace of $\F_2^Q$ spanned by $e_1,\ldots,e_{s(N)}$ where $e_i$ is the $i$-th
elementary basis vector.
For example, if $\hat{\tvar}_w = (\Read,\trole)$, then syntactically the
player's answer is a triple $(y,y^\perp,a)$ in $\F_2^Q \times \F_2^Q
\times \{0,1\}^*$.
We assume that the decider $\hat{\decider}^\intro$ computes the dimension $s(N)$ of the
subspace $V$ by calling $\sampler$ on input $(N,\gamestyle{dimension})$, and
 if $y,y^\perp$ are not presented as vectors in the subspace $V$,
then the decider rejects. Thus in the analysis we directly consider $y,y^\perp$ as vectors in $V$.
In \Cref{fig:intro-decider}, the components of the players' answers are
subscripted by the player index.
For example, if player $w$ receives question $(\Hide{k},\trole)$, then their
answers are denoted by $(y_w,y_w^\perp,x_w)$.

The notation used in the ``answer key'' is meant to give an indication of the
intended meaning of the players' answers.
We use $y$ to denote a vector that is supposed to be the result of measuring a
CL function $L^\trole$; $y^\perp$ is supposed to be the result of measuring
``dual'' linear maps $L^\perp$ (as used in Step~\ref{enu:hiding} in
Figure~\ref{fig:intro-decider}); $x$ is supposed to be the result of $\sigma^X$
measurements, and $z$ is supposed to be the result of $\sigma^Z$ measurements.
We use $a$ to denote the introspected answers meant for the original decider
$\decider$.

\medskip
\emph{CL functions and factor spaces.}
For $\trole \in \AB$, let $L^\trole = L^{\trole,\, N}$ denote the CL function
for original player $\trole$ specified by $\sampler$ on index $N = 2^n$.
For $z \in V$, the decider $\hat{\decider}^\intro$ computes $L^\trole(z)$ by calling
$\sampler$ on input $(N, \trole, \gamestyle{marginal}, \ell, z)$.

For $y \in V$ and $k \in \{1,\ldots,\ell\}$ we define register subspaces
$V_{k}^\trole (y)$ by induction on $k$.
For $k=1$, $V_1^\trole(y)$ is the first factor space\footnote{See
  Lemma~\ref{lem:cl-kth} for the definitions of factor spaces of a CL function.}
of $L^\trole$ and is independent of $y$.
Suppose $V_j^\trole(y)$ has been defined for all $j < k$.
Then we define the marginal space $V_{< k}^\trole(y) = \bigoplus_{j=1}^{k-1}
V_k^\trole(y)$, and define $V_k^\trole(y)$ as the $k$-th factor space $V_{k,\,
  u}^\trole$ of $L^\trole$ with prefix $u = L_{<k}^\trole(y)$.
We also define $V^\trole_{\leq k}(y) = V^\trole_{< k+1}(y)$, and $V^\trole_{>
  k}(y)$ to be the complementary register subspace to $V^\trole_{\leq k}(y)$
within $V$.

The decider $\tdint$ computes factor spaces $V_j^\trole(y)$ from $y \in
V$ in the following sequential manner: first, the indicator vector for the
factor space $V^\trole_1$ is computed by calling the original sampler $\sampler$
on input $(N,\trole,\gamestyle{factor},1,0)$.
Let $y_1$ denote $L_1^\trole(y)$.
Then, for $j \in \{2,\ldots,\ell\}$, the factor space $V^\trole_j(y)$ is
computed by calling $\sampler$ on input $(N,\trole,\gamestyle{factor},j,y_{\le
  j-1})$, where $y_{\le j-1} = L^\trole_{\leq j-1}(y)$.

Finally, \Cref{lem:cl-kth} implies that the CL function $L^\trole$ gives rise to functions $\{ L_{k,u}^\trole \}$ where for all $k$ and prefixes $u$, $L_{k,u}^\trole$ is a map from $V_{k,u}^\trole$ to $V_{k,u}^\trole$.

\medskip
\emph{Detailed explanation of the steps of $\tdint$.}
We give more details on the implementation of decider $\tdint$
specified in Figure~\ref{fig:intro-decider}.

\begin{enumerate}
\item The decider $\tdint$ first checks that the answers are not too
  long; \hnote{added some clarifications about answer-length check:} the maximum-length answer should be a tuple $(y, y^\perp, a)$ in response to question $(\Read,v)$ 
  where $y, y^\perp \in V$ and $a$ is an answer intended for the original
  decider $\decider$ on index $N$ (which we assume runs in time at most
  $N^\lambda$), or a tuple $(y,y^\perp,x) \in V \times V \times V$ in response to question $(\Hide{k},v)$. Either way, the maximum answer length should be $3Q = 3 \cdot 2^m \cdot \log q = \poly(R)$ bits long. 
  This check is necessary in order to ensure that the decider $\tdint$
  halts in time that is a fixed polynomial of $R = N^\lambda$.


\item If the question types for the players are both in $\type^\pauli$,
  $\tdint$ executes the decision procedure $\decider^\pauli$ for the
  Pauli basis test parametrized by $\introparams(R) = (q,m,d)$ (see
  Section~\ref{sec:qld-game} for the definition of $\decider^\pauli$).
  \hnote{edited, in response to Bowen comments} The Pauli basis test is intended to ensure that the players share at least $Q \geq R$ EPR pairs, where $Q = 2^m \log q$. This number of EPR pairs is chosen
  so that, under the assumption that the original verifier $\verifier$ is $\lambda$-bounded (see \Cref{def:lambda}), 
  the time complexity of $\sampler$ on index $N$ (and therefore its dimension) is at most $R = N^\lambda$. 
	When analyzing the Completeness, Soundness, and Entanglement properties of the introspective verifier $\verifier^\intro$,
	we assume that the original verifier $\verifier$ is $\lambda$-bounded (see \Cref{thm:introspection}). 
   


\item In Step~\ref{enu:sampling-pauli} of $\tdint$, player
  $w\in\{\alice,\bob\}$ receives question $(\Pauli,Z)$ and player $\overline{w}$
  receives question $(\Sample,\trole)$, for some $v\in\{\alice,\bob\}$.
  According to the answer key, player $w$ is expected to return an answer
  $\hat{a}_w\in \F_2^Q$ and player $\overline{w}$ is expected to respond with a
  pair $(y_{\overline{w}},a_{\overline{w}}) \in V \times \{0,1\}^*$.
  Thus, the dimension of answer $\hat{a}_w$ may be larger than that of
  $y_{\overline{w}}$; this is the reason that Step~\ref{enu:sampling-pauli}
  checks consistency between $y_{\overline{w}}$ and the projection of
  $\hat{a}_w$ to $V$.



\item In Step~\ref{enu:sampling-intro} of $\tdint$, player $w \in \ab$
  receives question $(\Introspect,\trole)$ and player $\overline{w}$ receives
  question $(\Sample,\trole)$.
  As specified by the ``answer key'', player $w$ responds with $(y_w,a_w) \in V
  \times \{0,1\}^*$ and player $\overline{w}$ responds with
  $(z_{\overline{w}},a_{\overline{w}}) \in V \times \{0,1\}^*$.
  The decider $\tdint$ checks that $a_w = a_{\overline{w}}$ and $y_w =
  L^\trole(z_{\overline{w}})$ where $L^\trole$ denotes the CL function for
  player $\trole$ specified by $\sampler$ for index $N = 2^n$.
  The CL function is computed by calling $\sampler$ on input
  $(N, \trole, \gamestyle{marginal}, \ell, z)$.

\item In Step~\ref{enu:hiding-read}, $\tdint$ checks that the answer of
  player $w$ who received $(\Hide{\ell},\trole)$ is consistent with the answer
  of player $\overline{w}$ who received $(\Read,\trole)$. One of the checks is
  that $y_{w,< \ell} = y_{\overline{w},<\ell}$~; these are, respectively, 
  $L^\trole_{< \ell}(y_w)$ and $L^\trole_{< \ell}(y_{\overline{w}})$.
  


\item\label{enu:decider-perp}
  \def\Lfunc{L^{\trole}_{k+1,\, y_{\overline{w},\, \le k}}}

  In Steps~\ref{enu:hiding-same}, the vectors $x_{w,\, >k+1}$ and
  $x_{\overline{w},\, >k+1}$ denote the projections of $x_w$ and
  $x_{\overline{w}}$ to the factor spaces $V^\trole_{> k+1}(y_w)$ and $V^\trole_{>
    k+1}(y_{\overline{w}})$, respectively.
  Similarly, $y_{\overline{w},\, k+1}^\perp$ denotes the projection of
  $y_{\overline{w}}^\perp$ to the factor space $V^\trole_{k+1} (y_{\overline{w}})$ and
  $x_{w,\, k+1}$ denotes the projection of $x_{w}$ to
 the factor space $V^\trole_{k+1}(y_{\overline{w}})$. We emphasize that the latter
 factor space depends on $y_{\overline{w}}$ and not $y_w$.\footnote{ This is because it is the $\overline{w}$ player
 who receives the $(\Hide{k+1},\trole)$ question, and in the ``honest'' strategy (described in \Cref{sec:intro-completeness}) 
 they are supposed to measure $k$ registers to obtain a sequence of values $(y_{\overline{w},1},y_{\overline{w},2},\ldots,y_{\overline{w},k})$ 
 which specify the $(k+1)$-st factor subspace, whereas the $w$ player is only supposed to measure $k-1$ registers.
 }


  The decider also has to compute $( \Lfunc )^\perp (x_{w,\, k+1})$.
  According to Definition~\ref{def:Lperp}, this requires specifying a basis for
  $\ker( \Lfunc )^\perp$.
  To compute the value, the decider performs the following steps:
	\begin{enumerate}
  \item Call $\sampler$ on input $(N, \trole, \gamestyle{factor}, k+1, y_{\overline{w},\, \le k})$ to
    obtain a subset $H = \{h_1,\ldots,h_m\}$ of the canonical basis for
    $\F_2^{s(N)}$ that is a basis of the register subspace $V^\trole_{k+1}(y_{\overline{w}})$.
  \item For $i \in \{1, 2, \ldots, m\}$ compute $c_i = \sampler(N, \trole,
    \gamestyle{linear}, k+1, y_{\overline{w},\, \le k}, h_i)$.
    Compute a matrix representation $M$ for $ \Lfunc $ in the basis
    $H$, whose columns are the vectors $c_1, \ldots, c_m$ as elements of
    $V^\trole_{k+1}(y_{\overline{w}})$.

  \item Using a canonical algorithm for Gaussian elimination, compute a basis
    $F$ for $\ker(M)$.

  \item Compute the canonical complement $S$ of $F$, as in
    Definition~\ref{def:canonical-complement}.
    $S$ is a basis for $\ker( \Lfunc )^\perp$.

  \item To compute $(\Lfunc)^\perp$ on input
    $x_{w,\, k+1}$, compute the canonical linear map with kernel basis
    $S$ (see Definition~\ref{def:cl-canonical}) on input $x_{w,\,
      k+1}$.
  \end{enumerate}

\item In Step~\ref{enu:hiding-pauli}, the player $w$ that receives $(\Pauli,X)$
  is expected to return an answer $\hat{a}_w$ in $\F_2^Q$.
  Part of this step checks that the projection of $\hat{a}_w$ to $V_{>1}$ is
  equal to $x_{\overline{w},\, >1}$ (which is the projection to $V_{>1}$ of the
  third component of the answer triple of player $\overline{w}$ that receives
  question $(\Hide{1},\trole)$).
\end{enumerate}

\subsubsection{Complexity of the introspective verifier}
\label{sec:intro-complexity}

In this section we determine the complexity of the introspective verifier.
The following lemma establishes the complexity of the decider $\tdint$
as well as the complexity of computing a {description} of it from the parameters
$(\lambda,\ell)$ and the description of the original verifier $\verifier =
(\sampler,\decider)$.

\begin{lemma}
  \label{lem:intro-decider-complexity}
  There exists a polynomial time Turing machine $\ComputeIntroDecider$ that
  on input $(\verifier, \lambda, \ell)$ outputs a description of a Turing machine where:
  \begin{enumerate}
  		\item   \hnote{jan 27, 2022: added this item} Its description length is at most $\poly(\lambda, \ell)$.
		\item Its time complexity is at most $\TIME_{\tdint}(n) = \poly(2^{\lambda n} , \ell)$.
  \end{enumerate}
  Furthermore, if $|\verifier| \leq \lambda$, then the output of $\ComputeIntroDecider$ is the introspective decider 
  $\tdint$ corresponding to verifier $\verifier$ and parameters $(\lambda,\ell)$. 
\end{lemma}

\begin{proof}
  Define the following $10$-input Turing machine $\RawIntroDecider$, \hnote{jan 27, 2022: added} which is 
  a universal Turing machine that computes the same function as the introspection decider $\tdint$, except
  it also takes as input extra inputs to specify the sampler, decider, and the parameters $\lambda, \ell$ that are 
  used in $\tdint$. Since $\RawIntroDecider$ is a universal Turing machine, its description length is constant.
  More formally, on input
  \begin{equation*}
    (\verifier', \lambda', \ell', n', x_1', x_2', \ldots, x_6'),
  \end{equation*}
    \hnote{jan 27, 2022, edited:} the Turing machine $\RawIntroDecider$ computes the output of the decider $\tdint$
    corresponding to $\verifier', \lambda', \ell'$ with input tapes set to $(n',x_1',\ldots,x_6')$.
In more detail, the Turing machine $\hat{\cal{D}}^{\intro}$ first computes
  $\introparams(R)$ for $R=2^{\lambda' n'}$. 
  Using Lemma~\ref{lem:introparams-complexity}, this computation takes time
  \hnote{edited in response to LB comment} $\polylog(R)$.
  It then executes the tests described in Figure~\ref{fig:intro-decider}.
  Write $\verifier' = (\sampler', \decider')$.
  The complexity of performing the entire procedure is subsumed by the
  complexity of the following steps:
	\begin{enumerate}
  \item Running the decider $\decider^\pauli$, which takes time $\poly(R)$ by
    Lemma~\ref{lem:qld-complexity};
  \item Running the decider $\decider'$ (on index $N' = 2^{n'}$) for at most
    $2^{\lambda' n'}$ steps;
  \item Running the sampler $\sampler'$ (on index $N'$) in order to
    compute the dimension $s(N')$ and the marginal and factor spaces, and the CL
    functions as described in Section~\ref{sec:intro-verifier}.
    $\sampler'$ is called at most $\poly(s(N'),\ell')$ times; due to the timeout,
    each call takes time at most $2^{\lambda' n'}$.
  \item Computing $(L^\trole_{k,\, u})^\perp(x_{\overline{w},\, k})$ in
    Step~\ref{enu:hiding-same}.
    This only requires to perform Gaussian elimination and other simple finite
    field manipulations that can be implemented in $\poly(s(N'))$
    time.
	\end{enumerate}
  All other tests are elementary.
  Thus the time complexity of $\RawIntroDecider$ is $\poly(R,\ell')$.
  Note that the bound is independent of $\verifier'$: due to the abort
  condition in the definition of $\tdint$, the Turing machine
  $\hat{\decider}^\intro$ aborts if the runtime of $\sampler'$ or $\decider'$ is
  larger than $2^{\lambda' n'}$. 

  \hnote{jan 27, 2022: edited this part: }We now define the Turing machine $\ComputeIntroDecider$: on input
  $(\verifier,\lambda,\ell)$, it first determines whether the description length
	$|\verifier|$ of the verifier is at most $\lambda$.\footnote{This can be done by keeping a counter of $O(\log \lambda)$ bits, and incrementing it as the input $\verifier$ is being scanned through. If at any point the counter exceeds $\lambda$, the scanning stops.} If it isn't, then it sets $\verifier' = (\sampler',\decider')$ 
	where $\sampler',\decider'$ are the trivial Turing machine whose description is $0$. Otherwise,
  $\ComputeIntroDecider$ sets $\verifier' = \verifier$. Note that by construction, the description length of $\verifier'$ is at most $\lambda$, 
  and that the time complexity of computing the description of $\verifier'$ is polynomial in $|\verifier|$ and $\log \lambda$.
  
  Then $\ComputeIntroDecider$ outputs the description of
  $\RawIntroDecider$ with the first three input tapes hardwired to
  $\verifier'$, $\lambda$, and $\ell$, respectively, yielding the decider
  $\tdint$ corresponding to the verifier $\verifier'$ and
  parameters $(\lambda, \ell)$. 
  Computing this description takes $\poly(\abs{\verifier}, \log\lambda, \log
  \ell)$ time. Thus overall the complexity of $\ComputeIntroDecider$ is $\poly(\abs{\verifier},\log \lambda,\log \ell)$, which is polynomial in its input length.
  
  \hnote{jan 27, 2022: added} The description length of $\tdint$ is simply the description length of $\RawIntroDecider$ (which is a universal constant),
  plus the description length of $\verifier'$ (which is at most $\lambda$), plus the description lengths of $\lambda$ and $\ell$, which altogether
  totals $\poly(\lambda, \ell)$ as desired. The time complexity of $\tdint$ follows from the time complexity of $\RawIntroDecider$, which we established is $\poly(R, \ell) = \poly(2^{\lambda n},\ell)$.
  
  Finally, we note that if $|\verifier| \leq \lambda$, then by construction the output of $\ComputeIntroDecider$ is the decider $\tdint$ corresponding to the verifier $\verifier$ and parameters $(\lambda,\ell)$.
\end{proof}

\subsubsection{Introspection theorem}
\label{sec:intro-theorem}

We are now ready to state the introspection theorem.

\begin{theorem}[Introspection theorem]
  \label{thm:introspection}
  There is a polynomial time Turing machine $\ComputeIntroVerifier$ which takes
  as input a tuple $(\verifier, \lambda, \ell)$ for $\lambda, \ell \in \N$ and
  returns the description of the introspective verifier $\vint$.
  For all $\ell \in \N$, $\vint$ is a $5$-level normal-form verifier
  with complexity measures
  \begin{equation*}
    \begin{split}
      \TIME_{\sint}(n) & = \poly \big (n, \lambda, \ell \big),\\
      \TIME_{\dint}(n) & = \poly \big (2^{\lambda n}, \ell \big),\\
      |\dint| &= \poly(\lambda,\ell).
    \end{split}
  \end{equation*}
  Moreover, for all $\ell$, there exist constants $a>0$, \, $0 < b < 1$ (depending on $\ell$),
  and a function $\delta(\eps, n) = a \bigl( (\lambda n)^a\eps^b + (\lambda
  n)^{-b} \bigr)$ such that for all $n \in \N$ and $\eps \ge 0$ the following
  statements hold if $\verifier$ is a $\lambda$-bounded $\ell$-level verifier.
  \begin{enumerate}
  \item (\textbf{Completeness}) If $\verifier_{2^n}$ has a projective, consistent,
    and commuting (PCC) strategy with value $1$ then $\vint_n$ also
    has a PCC strategy with value $1$.

  \item (\textbf{Soundness}) If $\val^*(\vint_n) > 1 - \eps$, then
    $\val^*(\verifier_{2^n}) \geq 1 - \delta(\eps,n)$.
  \item (\textbf{Entanglement}) The entanglement bound $\Ent$ defined
    in \Cref{def:ent} satisfies
    \begin{equation*}
      \Ent(\vint_n, 1 - \eps) \geq \max
      \left \{ \Ent(\verifier_{2^n}, 1 - \delta(\eps,n)),
        (1 - \delta(\eps,n))\, 2^{2^{\lambda n}} \right \}\;.
    \end{equation*}
  \end{enumerate}
\end{theorem}

\begin{proof}
  The Turing machine $\ComputeIntroVerifier$ does the following on input
  $(\verifier, \ell)$: it first computes
  \begin{align*}
    \tsint & = \ComputeIntroSampler(\lambda,\ell)\;,\\
    \tdint & = \ComputeIntroDecider(\verifier, \lambda, \ell)\;,
  \end{align*}
  using \cref{lem:intro-sampler-complexity,lem:intro-decider-complexity}, runs
  the detyping procedure from \cref{def:detyped-verifier} to compute
  \begin{equation*}
    \vint = \detype(\tvint)
  \end{equation*}
  where $\tvint = (\tsint,\tdint)$ and then outputs the resulting detyped verifier.
  This takes time $\poly(\abs{\verifier}, \log\lambda, \log\ell)$, which is polynomial
  in the input length of $\ComputeIntroVerifier$. 
  The complexity parameters of the typed sampler $\tsint$ and typed decider
  $\tdint$ follow from
  \cref{lem:intro-sampler-complexity,lem:intro-decider-complexity}.
  The complexity parameters of the detyped verifier $\vint$ then follow from
  \cref{lem:detyping-verifiers}.
  As $\tvint$ is a $3$-level verifier, $\vint$ is $5$-level as detyping will
  increase the number of levels by at most $2$.

  The completeness part of the theorem is proven in \cref{sec:intro-completeness} and
  the soundness and entanglement properties are proven in \cref{sec:intro-soundness}.
\end{proof}


\subsection{Completeness of the introspective verifier}
\label{sec:intro-completeness}

In this section we establish the completeness property of the introspection
game: if for $N = 2^n$, $\verifier_N$ has a PCC strategy (see
Definition~\ref{def:spcc}) with value $1$, then so does $\vint_n$.
For this we describe the actions that are expected of the players in the
introspection game (i.e.\ the ``honest strategy'').
We first prove several preliminary lemmas that will be used in both the
completeness and soundness analysis.

\subsubsection{Preliminary lemmas}

The lemmas in this section are stated for general fields $\F$ and generalized
Pauli observables and projectors, although in the application to introspection
games we use $\F = \F_2$, $\omega = -1$, and qubit Pauli observables and
projectors.
Furthermore, the Pauli observables $\tau^W(v)$ and projectors $\tau^W_u$ act on
$\C^{\F^k}$ for some integer $k$ (in our application, we set $k=R$).



The following lemma generalizes Eq.~\eqref{eq:pauli-inversion-0}.

\begin{lemma}\label{lem:why-didnt-i-think-of-this-before}
Let $L: \F^k \to \F^k$ be a linear map, and let $W \in \{X, Z\}$.
For each $a$ in the range of $L$, let $u_a \in \F^k$ be such that $L(u_a) = a$. 
\begin{enumerate}
\item For each $v \in \ker(L)^\perp$, 
\begin{equation*}
\tau^W(v) = \sum_{a \in \F^k} \omega^{\tr(u_a \cdot v)} \tau^W_{[L(\cdot)=a]}\;.
\end{equation*}
\item For all $a$ in the range of $L$,
\begin{equation*}
\tau^W_{[L(\cdot)=a]} = \E_{v \sim \ker(L)^\perp} \omega^{-\tr(v \cdot u_a)} \tau^W(v)\;.
\end{equation*}
\end{enumerate}
\end{lemma}

\begin{proof}
  Let $V$ denote the image of $\F^k$ under $L$.
  Let $a\in V$, $v \in \ker(L)^\perp$.
  For all $u, u' \in L^{-1}(a)$, we have that $u-u'\in \ker(L)$ and thus $\tr(u
  \cdot v) = \tr(u' \cdot v)$.
  As a result, using~\eqref{eq:pauli-obs-proj}, for any $v\in\ker(L)^\perp$ it
  holds that
  \begin{equation*}
    \tau^W(v)
    = \sum_{u \in \F^k} \omega^{\tr(u \cdot v)} \tau^W_u
    = \sum_{a \in V} \sum_{u \in L^{-1}(a)} \omega^{\tr(u \cdot v)} \tau^W_u
    = \sum_{a \in V} \omega^{\tr(u_a \cdot v)} \tau^W_{[L(\cdot)=a]}
    = \sum_{a \in \F^k} \omega^{\tr(u_a \cdot v)} \tau^W_{[L(\cdot)=a]}\;,
  \end{equation*}
  where in the last equality we used that for $a \notin V$, the projector
  $\tau^W_{[L(\cdot)=a]}$ vanishes.
  This shows the first item.
  For the second item,
  \begin{align*}
    \tau^W_{[L(\cdot)=a]}
    & = \sum_{u  \in L^{-1}(a)} \tau^W_u\\
    & =  \sum_{u  \in L^{-1}(a)} \E_{v \sim \F^k} \Bigl(
      \omega^{-\tr(u \cdot v)} \tau^W(v) \Bigr) && \\
    & =  \sum_{u \in \ker(L)} \E_{v \sim \F^k} \Bigl(
      \omega^{-\tr((u_a + u) \cdot v)} \tau^W(v) \Bigr)\\
    & = \frac{|\ker(L)|}{|\F^k|} \sum_{v \in \F^k} \biggl( \Big(
      \E_{u \sim \ker(L)} \omega^{-\tr(u \cdot v)} \Big)\,
      \omega^{-\tr(u_a \cdot v)}  \tau^W(v) \biggr)\\
    & = \E_{v \in \ker(L)^\perp} \omega^{-\tr(v \cdot u_a)} \tau^W(v)\;,
\end{align*}
where the second equality follows from~\eqref{eq:pauli-inversion-0} and the last
uses the fact that $|\F^k|=|\ker(L)||\ker(L)^\perp|$, as shown in
Lemma~\ref{lem:perp_perp}, and Lemma~\ref{lem:cancellation} applied to the
expectation over~$u$.
\end{proof}


\begin{lemma}[Commuting~$X$ and~$Z$ measurements]
  \label{lem:commute} 
	For all linear maps $L,R: \F^k \to \F^k$ such that
  \begin{equation*}
    \ker (R)^\perp \subseteq \ker (L)\;,
  \end{equation*}
  the measurements
  $\big\{ \tau^Z_{[L(\cdot)=b]} \big\}_{b\in \F^k}$ and $\big\{ \tau^X_{[R(\cdot) = d]}
  \big\}_{d\in \F^k}$ commute.
\end{lemma}
\begin{proof}
  Let $b,d\in \F_k$.
  If either $b$ is not in the range of $L$, or $d$ is not in the range of $R$,
  then at least one of $\tau^Z_{[L(\cdot) = b]}$ or $\tau^X_{[R(\cdot) = d]}$ is
  $0$, and thus the operators trivially commute.
  Otherwise, both $b$ and $d$ are in the range of $L$ and $R$, respectively.
  Let $a_0 \in L^{-1}(b)$, $c_0 \in R^{-1}(d)$.
  By Lemma~\ref{lem:why-didnt-i-think-of-this-before},
  \begin{equation*}
    \tau^Z_{[L(\cdot)=b]} = \E_{u \in \ker(L)^\perp}
    \omega^{-\tr(u \cdot a_0)} \tau^Z(u)\;,
  \end{equation*}
  \begin{equation*}
    \tau^X_{[R(\cdot)=d]} = \E_{v \in \ker(R)^\perp}
    \omega^{-\tr(v \cdot c_0)} \tau^X(v)\;.
  \end{equation*}
  For any $v \in \ker(R)^\perp$, by assumption $v\in \ker(L)$, so for $u \in
  \ker(L)^\perp$ it holds that $u\cdot v=0$.
  Thus $\tau^Z(u)$ and $\tau^X(v)$ commute, and it follows that
  $\tau^Z_{[L(\cdot)=b]}$ and $\tau^X_{[R(\cdot)=d]}$ commute as well.
\end{proof}

\subsubsection{Completeness of the introspective verifier}
\label{sec:intro-completeness-sub}













\begin{proof}[Proof of the completeness part of \cref{thm:introspection}]
  We analyze the completeness of the \emph{typed}
  verifier $\tvint$; the corresponding completeness of the
  \emph{detyped} verifier $\vint = \detype(\tvint)$ follows from
  Lemma~\ref{lem:detyping-verifiers}, and the fact that the type set
  $\type^\intro$ has size $O(\ell)$.

  \paragraph{Completeness.}
  We first show completeness.
  Let $n\geq 1$ be an index for $\tvint$ and $N = 2^n$ be the corresponding
  index for $\verifier$.
  Since we assume that the verifier $\verifier$ is $\lambda$-bounded, we have
  that
  \begin{equation}
    \label{eq:intro-complexity-assump}
    \max \big\{ \TIME_\sampler(N), \, \TIME_\decider(N) \big\} \leq N^\lambda = R\;.
  \end{equation}
  This assumption on the time complexity of $\verifier$ ensures that
  $\tdint$ never aborts due to a timeout.
	Let $L^\trole = L^{\trole,\, N}$ denote the CL function of the original
  sampler $\sampler$ on index $N$ corresponding to player $\trole \in \AB$.
  Let $R = N^\lambda$, and let $\introparams(R) =
  (q,m,d)$, as in Section~\ref{sec:intro-verifier}.
  Set $Q = M\log q$ where $M=2^m$\,; this represents the number of qubits that are
  certified by the Pauli basis test parameterized by $\introparams(R)$.
  
  Let $\strategy = (\ket{\aux}, A, B)$ be a PCC strategy for $\verifier_N$ with
  value~$1$.
  We first construct a PCC strategy $\strategy^\intro_n$ for the typed verifier
  $\tvint_n$ with value $1$.
  We then conclude using Lemma~\ref{lem:detyping-verifiers}.

  \begin{figure}[htb!]
    \centering
    \begin{gamespec}
      \newcolumntype{M}{>{\arraybackslash $}X<{$}}
      \renewcommand{\arraystretch}{1.6}
      \begin{tabularx}{\textwidth}{*{6}M}
        \toprule
        & \phantom{I1} V^\trole_1
        & \phantom{I1} V^\trole_2 & \phantom{I1} V^\trole_3 & \aux \\
        \midrule
        (\Introspect, \trole) & \sigma^Z_{L_1} & \sigma^Z_{L_2} &
        \sigma^Z_{L_3} & A^x/B^x \\
        (\Sample, \trole) & \sigma^Z & \sigma^Z & \sigma^Z & A^x/B^x \\
        (\Read, \trole) & \sigma^Z_{L_1^{\phantom{\perp}}}\;
        \sigma^X_{L_1^\perp} &
        \sigma^Z_{L_2^{\phantom{\perp}}}\;
        \sigma^X_{L_2^\perp} &
        \sigma^Z_{L_3^{\phantom{\perp}}}\;
        \sigma^X_{L_3^\perp} & A^x/B^x \\
        \midrule
        (\Hide{3}, \trole) & \sigma^Z_{L_1^{\phantom{\perp}}}\;
        \sigma^X_{L_1^\perp} &
        \sigma^Z_{L_2^{\phantom{\perp}}}\;
        \sigma^X_{L_2^\perp} &
        \phantom{\sigma^Z_{L_3^{\phantom{\perp}}}\;}
        \sigma^X_{L_3^\perp} & I \\
        (\Hide{2}, \trole) & \sigma^Z_{L_1^{\phantom{\perp}}}\;
        \sigma^X_{L_1^\perp}
        & \phantom{\sigma^Z_{L_1^{\phantom{\perp}}}\;}
        \sigma^X_{L_2^\perp}
        & \sigma^X & I \\
        (\Hide{1}, \trole) & \phantom{\sigma^Z_{L_1^{\phantom{\perp}}}\;}
        \sigma^X_{L_1^\perp} & \sigma^X & \sigma^X & I \\
        \midrule
      \end{tabularx}
      The left-most column denotes the introspection/hiding questions that a
      player may receive.
      The top row denotes the registers corresponding to the factor spaces of a
      CL function $L^\trole$ (we note that the partition of the registers depends
      on the prefixes), as well as the register corresponding to the state
      $\ket{\aux}$ coming from the original PCC strategy $\strategy$.
      We use $\sigma^Z_{L_j}$ as shorthand for $\sigma^Z_{[ L^\trole_{j,\,
          x_{<j}} (\cdot) = x_j ]}$ and similarly $\sigma^X_{L^\perp_j}$ for $
      \sigma^X_{[ (L^\trole_{j,\, x_{<j}})^\perp (\cdot)=x_j^\perp]}$.
      A symbol $I$ means that the register is left unmeasured.
    \end{gamespec}
    \caption{Summary of the honest strategy $\strategy_n^\intro$ for
      $\hat{\verifier}_n^\intro$ for a $3$-level sampler.}
    \label{fig:intro-honest}
  \end{figure}

  \begin{remark}
    \label{rk:downsizing}
    Note that by definition in $\verifier_n^\intro$ the players receive
    questions $(x,y)$ that are sampled according to the distribution
    $\mu_{\tsint,\, n}^{G^\intro}$ associated with the downsized typed
    sampler $\downsize(\tilde{\sampler}^\intro)$.
    Using the definition of the downsized typed sampler,
    Definition~\ref{def:downsize-typed-sampler}, and
    Lemma~\ref{lem:downsize-cl-dist} the distribution is identical to the
    distribution $\mu_{\tilde{\sampler}^\intro,\, n}^{G^\intro}$, up to the
    bijective mapping $\downsize$.
    This mapping can be computed by the players themselves.
    Therefore, we construct a strategy for players that receive questions from
    $\mu_{\tilde{\sampler}^\intro,\, n}^{G^\intro}$, and a strategy for
    questions from $\mu_{\tsint,\, n}^{G^\intro}$ follows immediately.
  \end{remark}

  The strategy $\strategy^\intro_n$ uses the state $\ket{\EPR_2}^{\otimes (Q+1)}
  \otimes \ket{\aux}$, where recall that $\ket{\EPR_2} =
  \frac{1}{\sqrt{2}}(\ket{00}+\ket{11})$.
  For all register subspaces $K \subseteq \F_2^{Q}$ (see
  Definition~\ref{def:cl-func} for the definition of register subspace),
  whenever we refer to ``register $K$'', we mean the qubits of
  $\ket{\EPR_2}^{\otimes Q}$ corresponding to $K$ (see
  Section~\ref{sec:lin-reg}).
  The most frequent register subspace we consider is~$V$, spanned by $e_1,
  \ldots, e_{s(N)}$.
  We write $\overline{V}$ for the complement of~$V$, i.e.
  the register subspace spanned by $e_{s(N)+1}, \ldots, e_{Q}$.
  (Note that $s(N) \leq R$ by
  assumption~\eqref{eq:intro-complexity-assump}, \hnote{added:} and by \Cref{lem:delta-bound} we have $R \leq M \leq Q$.)
  Then $\ket{\EPR_2}^{\otimes Q} = \ket{\EPR}_V \otimes
  \ket{\EPR}_{\overline{V}}$.

  Let $\strategy^\pauli_n$ be the honest Pauli strategy with respect to
  $\introparams$ defined in the proof of \cref{lem:pauli-completeness}. Using the 
  isomorphism between $\C^q$ and $(\C^2)^{\otimes \log q}$ specified by \Cref{lem:pauli-binary},
 the measurements   of $\strategy^\pauli_n$ can be treated as qubit Pauli measurements acting on the
  maximally entangled state over qubits. Furthermore the questions and answer labels
  of the measurements can be viewed as binary strings
 using the bijection $\kappa: \F_q \to \F_2^{\log q}$ specified in \Cref{sec:finite-fields}.
  For a question with a type from $\type^\pauli$ the player measures the shared state
  $\ket{\EPR_2}^{\otimes (Q+1)}$ using the measurements specified by
  $\strategy^\pauli_n$, and reports the measurement outcomes.
  When a player receives questions of type $\hat{\tvar} \in \type^\intro
  \setminus \type^\pauli$, they perform measurements described as follows.
  (Below, whenever we write a Pauli operator $\sigma^W_a$ the register on which
  the operator acts should always be clear from context, and is implicit from
  the space in which the outcome $a$ ranges.)
  The reader may find it helpful to consult Figure~\ref{fig:intro-honest} to see
  a summary of the honest strategy $\strategy^\intro_n$ for the special case
  when $\ell = 3$.

  \begin{description}
  \item $(\Introspect, \trole)$: The player performs the measurement
    \begin{equation}
      \label{eq:main-measurement}
      \bigl\{ \sigma^Z_{[L^\trole(\cdot) = y]} \bigr\}
    \end{equation}
    to obtain an $y\in V$.
    Intuitively, the player has now introspectively sampled the question $y$ for
    original player $\trole$ in game~$\verifier_N$.
    The player then measures $\ket{\aux}$ using player $\alice$'s
    measurement~$\{A^y_a\}$ from $\strategy$ if $v = \alice$ and using player
    $\bob$'s measurement~$\{B^y_a\}$ if $\trole = \bob$ to obtain an answer~$a$.
    The player replies with $(y,a)$.
  \item $(\Sample, \trole)$: The player measures their share of $\ket{\EPR}_V$
    in the $Z$ basis to obtain $z \in V$.
    Using this, they compute the question $y = L^\trole(z)$.
    The player then uses player $v$'s strategy and question $y$ to measure
    $\ket{\aux}$ and obtain outcome $a$.
    The player replies with $(z, a)$.
  \item $(\Read, \trole)$: The player first performs all measurements as in the
    $(\Introspect, \trole)$ question for player $\trole$ and records the
    outcomes as $y\in L(V)$ and $a\in\{0,1\}^*$.
    For $j\in \{1, 2, \ldots, \ell\}$, the player measures their share of the state
    $\ket{\EPR}_V$ with the measurement
\begin{equation}
      \label{eq:perp-measurement}
      \big\{\sigma^X_{[L_j^\perp (\cdot) = y_j^\perp]} \big\}_{y_j^\perp}
    \end{equation}
    to obtain $y^\perp = y_1^\perp + \cdots + y_\ell^\perp$.
    Here $L_j^\perp$ is shorthand for the function $(L^{\trole}_{j,\,
      y_{<j}})^\perp$ defined in Item~\ref{enu:decider-perp} of the decider
    description in Section~\ref{sec:intro-verifier}.
    (That these operators are simultaneously measurable with the measurement
    in~\eqref{eq:main-measurement} follows from Lemma~\ref{lem:commute} and the
    fact that $\ker(L^\perp_{j})^\perp = \ker(L_{j})$ by
    Lemma~\ref{lem:perp_perp} and the definition of $L^\perp_{j}$ in
    Section~\ref{sec:intro-verifier}.)
    
    The player measures $\ket{\aux}$ with player $\trole$'s measurement strategy
    in $\strategy$ for question $y$ to obtain $a$ and replies with
    $(y,y^\perp,a)$.

  \item $(\Hide{k}, \trole)$:
    The player performs the following sequence of measurements: first measure
    $\{ \sigma^Z_{[L^\trole_{1}(\cdot) = y_{1}]} \}$ on register $V_1^v$ to
    obtain $y_1$.
    Then, use $y_1$ to specify the second linear function
    $L^\trole_{2,\, y_1}(\cdot)$ and measure register $V^\trole_{2,\,y_1}$ using
    $\{ \sigma^Z_{[L^\trole_{2,\, y_1}(\cdot) = y_{2}]} \}$ to obtain $y_2$.
    This process continues until the $(k-1)$-th linear map $L_{k-1,\, y_{<
        k-1}}^\trole(\cdot)$ has been measured to obtain $y_{k-1}$ in factor
    space $V_{k-1,\, y_{< k-1}}^\trole$.
    Let $y = y_1 + y_2 + \cdots + y_{k-1}$.
    
    Next, for $j \in \{1, 2, \ldots, k\}$ the player measures
    \begin{equation*}
      \Big \{ \sigma^X_{[L_{j}^\perp(\cdot) = y^\perp_{j}]} \Big \}_{y^\perp_j}\;,
    \end{equation*}
    where $L_j^\perp$ denotes the linear map $(L^\trole_{j,\, y_{<j}})^\perp$ as in
    the case $(\Read, \trole)$.
    Let $y^\perp = y_1^\perp + y_2^\perp + \cdots + y_k^\perp$, where each
    $y_j^\perp$ is a vector in the factor space $V^\trole_{j,\, y_{<j}}$.
    Finally, the player measures register $V_{> k}^\trole(y)$ using
    $\{\sigma^X_{x_{>k}}\}$ to obtain outcome $x_{> k}$.
    Let $x = x_{> k}$.
    The player replies with $(y,y^\perp,x)$.
\end{description}

By definition, when player $w$ performs the honest measurement for question
$(\Introspect, \alice)$ and player $\overline{w}$ performs the honest
measurement for question $(\Introspect, \bob)$, the joint outcome $(y, y')$ has
distribution $\mu_{\sampler,\, N}$.
In this case, the players play according to strategy~$\strategy$ and succeed
with probability~$1$.
In all other cases, it is straightforward to verify that the players succeed in
all tests performed by $\tdint$ (Figure~\ref{fig:intro-decider}) with
probability~$1$.
As a result, the value of this strategy is~$1$. 

The strategy $\strategy^\intro$ is projective by construction.
It is also consistent because of the assumed consistency of the strategy
$\strategy$ as well as consistency of the honest Pauli strategy $\strategy^\pauli$.
Furthermore, note that $\strategy^\intro$ only calls $\strategy$ for both
players on question pairs such that both types are in
$\{\Introspect,\Sample,\Read\}$.
In all these cases, $\strategy$ is called on a pair of questions $(y,y')$
distributed as $(L^\trole(z),L^{\trole'}(z))$ for
$\trole,\trole'\in\{\alice,\bob\}$ and $z$ uniform in $V$.
When $v \neq v'$, any such pair by definition has positive probability under
$\mu_{\sampler,\, n}$, and so by assumption the associated measurements from
$\strategy$ commute.
On the other hand, when $v = v'$, then the players apply the same measurements
from $\strategy$, and because $\strategy$ only uses projective measurements,
their measurements commute as well.
Examining all other cases, it follows by direct inspection that the strategy
commutes on all question pairs whose corresponding types appear as an edge in
the graph $G^\intro$.
Thus the strategy commutes with respect to the support of the distribution
$\mu_{\tsint,\, n}$.
\end{proof}

\begin{remark}\label{rk:intro-normalform}
  For future reference, we note that on any input $(\verifier, \lambda, \ell)$,
the Turing machine $\ComputeIntroVerifier$ \emph{always} returns a normal form verifier $\tvint =
  (\tsint,\tdint)$ --- even if $\verifier$ itself is not the description of a normal form verifier. 
  This is because, for any two integer $\lambda,\ell$,
  $\ComputeIntroSampler(\lambda,\ell)$ returns a sampler with field size $2$,
  and for any $\verifier = (\sampler,\decider), \lambda$ and $\ell$ the decider $\tdint$ specified in
  Figure~\ref{fig:intro-decider} takes $7$ inputs and always halts with a
  single-bit output, even if $\sampler$ or $\decider$ themselves do not halt.
\end{remark}

\subsection{Soundness of the introspective verifier}
\label{sec:intro-soundness}

The main result of this section is the soundness part of
\cref{thm:introspection}.





\subsubsection{The Pauli twirl}

A key tool in the proof of soundness is the \emph{Pauli twirl}.
In this section we introduce the Pauli twirl and establish several of its
properties.
The section closely follows Sections~$8$ and~$10$ of~\cite{NW19}. 

To begin, we define the twirl with respect to an arbitrary distribution over
unitaries.

\begin{definition}[Twirl]
  \label{def:mixing}
  Let $\mu$ be a probability distribution over a finite set of unitary matrices.
  Then for any matrix $A$, the \emph{twirl of~$A$ with respect to~$\mu$},
  denoted $\mathscr{T}_\mu(A)$, is defined as
  \begin{equation*}
    \mathscr{T}_\mu (A) = \E_{U \sim \mu} \bigl( U A U^\dagger \bigr)\;.
  \end{equation*}
\end{definition}

In the next two lemmas we consider the Pauli twirl, in which the distribution
$\mu$ is over subsets of Pauli observables.
First, we show how the Pauli twirl acts on Pauli matrices.
Then, using this, we derive an expression for the Pauli twirl applied to general
matrices.

\begin{lemma}[Pauli twirl of Pauli matrices]
  \label{lem:twirl-pauli}
  Let $V$ be a subspace of $\F^k$, let $W \in \{X, Z\}$, and let $\mu$ be the
  uniform distribution over~$\{\tau^W(w) : w \in V\}$.
  Let $W'\neq W$ and $u\in \F^k$.
  Then
  \begin{equation*}
    \mathscr{T}_\mu(\tau^{W'}(u)) =
    \begin{cases}
      \tau^{W'}(u)&\text{ if }u \in V^{\perp}\;,\\
      0 &\text{ if }u \notin V^{\perp}\;.
    \end{cases}
  \end{equation*}
\end{lemma}

\begin{proof}
  Let $u\in \F^k$.
  Let $c = 1$ if $W = Z$ and let $c = -1$ if $W = X$.
  Then
  \begin{align*}
    \mathscr{T}_\mu(\tau^{W'}(u))
    &= \E_{z \sim V} (\tau^W(z)  \tau^{W'}(u) \tau^W(z)^\dagger)\\
    &= \E_{z \sim V} (\omega^{c\cdot \tr(u \cdot z)} \tau^{W'}(u) \tau^W(z)
      \tau^W(z)^\dagger)\\
    &= \Big(\E_{z \sim V} \omega^{c\cdot \tr(u \cdot z)}\Big) \tau^{W'}(u)\;.
  \end{align*}
  The lemma now follows from Lemma~\ref{lem:cancellation}
  and the fact that $c \cdot u \in V^\perp$ if and only if~$u\in V^\perp$.
\end{proof}

\begin{lemma}[Pauli twirl of general matrices]
  \label{lem:mixing-L}
  Let $V= \F^k$, and let $L : V \rightarrow V$ be a linear map.
  Let $\zeta$ be the uniform distribution over $\{\tau^Z(z) \mid z \in V\}$ and
  $\chi$ the uniform distribution over $\{\tau^X(x) \mid x \in \ker(L)\}$.
  Let~$M$ be a matrix acting on $\C^V \otimes \mH_\alice$, where $\mH_\alice$ is
  a finite dimensional Hilbert space.
  Then there exist matrices $\{M^y\}_{y\in L(V)}$ acting on $\mH_\alice$ such
  that the twirl of~$M$ with respect to~$\zeta$ and~$\chi$ is given by
  \begin{equation}
    \label{eq:mixing-L-0}
    (\mathscr{T}_{\chi} \circ \mathscr{T}_{\zeta} \otimes \Id_A)(M) = \sum_{y\in
      V} \tau^Z_{[L(\cdot)= y]} \otimes M^y\;.
  \end{equation}
	Moreover, if we apply~\eqref{eq:mixing-L-0} to each element of a POVM
  measurement $\{M_a\}$, then for each $y \in V$, the set $\{M^y_a\}$ also forms
  a POVM measurement.
\end{lemma}

\begin{proof}
  The collection $\{\tau^X(x)\tau^Z(z)\}_{x, z \in V}$ forms a basis for the
  complex linear space of matrices acting on $\C^V$.
  As a result, we can write
  \begin{equation*}
    M = \sum_{x, z \in V} \tau^X(x) \tau^Z(z) \otimes M_{x, z}\;,
  \end{equation*}
  for matrices $M_{x,z}$ on $\mH_\alice$.
  We now use Lemma~\ref{lem:twirl-pauli} to compute the twirl first with respect
  to $\zeta$ and then with respect to $\zeta$ and $\chi$:
  \begin{align*}
    (\mathscr{T}_{\zeta} \otimes \Id_{A})(M)
    &= \sum_{x, z \in V} \mathscr{T}_{\zeta}(\tau^X(x)) \tau^Z(z) \otimes M_{x, z}
      = \sum_{z \in V} \tau^Z(z) \otimes M_{0, z}\;,\\
    (\mathscr{T}_{\chi} \circ \mathscr{T}_{\zeta} \otimes \Id_{A})(M)
    &= \sum_{z \in V} \mathscr{T}_{\chi}(\tau^Z(z)) \otimes M_{0, z}
      =  \sum_{z \in \ker(L)^\perp} \tau^Z(z) \otimes M_{0, z}\;.
  \end{align*}
  For all $y \in V$, let $u_y$ denote an arbitrary element of $L^{-1}(y)$ if $y$
  is in the image of $L$; otherwise, set $u_y = 0$.
  Expanding $\tau^Z(z)$ using the first part of
  Lemma~\ref{lem:why-didnt-i-think-of-this-before},
  \begin{align}
    \label{eq:post-twirl-matrix}
    \sum_{z \in \ker(L)^\perp} \tau^Z(z) \otimes M_{0, z}
    & = \sum_{z \in \ker(L)^\perp} \,\, \sum_{y} \omega^{\tr(u_y \cdot z)}
      \tau^Z_{[L(\cdot)=y]} \otimes M_{0, z} \notag \\
    & =  \sum_{y} \tau^Z_{[L(\cdot)=y]} \otimes \Big(\sum_{z \in \ker(L)^\perp}
      \omega^{\tr(u_y \cdot z)}  M_{0, z}\Big)\;.
  \end{align}
  Equation~\eqref{eq:mixing-L-0} follows by setting $M^y = \sum_{z \in
    \ker(L)^\perp} \omega^{\tr(u_y \cdot z)} M_{0, z}$.

  For the ``moreover'' part, note first that whenever $M \geq 0$ it holds that
  any twirl satisfies $0 \leq \mathscr{T}_{\mu}(M)$.
  As a result, each matrix $M^y$ must be positive semi-definite due to
  Equation~\eqref{eq:post-twirl-matrix} and the fact that the
  $\{\tau^Z_{[L(\cdot)=y]}\}_y$ matrices are orthogonal projections.
  Next, suppose $\{M_a\}$ is a POVM measurement, and write $N = \sum_a M_a$ for
  the identity matrix.
  Then by linearity, for each $y \in V$, $\sum_a M_a^y = N^y$.
  In addition, $N_{0, 0} = \Id$, and $N_{x, z} = 0$ otherwise.
  As a result, $N^y = N_{0, 0} = \Id$, and so $\{M_a^y\}$ also forms a POVM.
\end{proof}

In the next few lemmas we derive a sufficient condition for a measurement to be
close to its own Pauli twirl, namely that it satisfies certain commutation
relations with the Pauli basis measurements.

\begin{lemma}[Commuting with Pauli basis implies commuting with Pauli observables]
  \label{lem:W} 
  Let~$\cal{X},\cal{A}$ be finite sets and~$D$ be a distribution over $\cal{X}$.
  For each~$x \in \cal{X}$, let $V_x$ be a register subspace of $V = \F^k$, and
  let $L_x:V_x \to V_x$ be a linear map.

  Consider a state $\ket{\psi} = \ket{\EPR}_V \otimes \ket{\aux}$, where
  $\ket{\EPR}_{V} \in \mH_\alice \otimes \mH_\bob$, for $\mH_\alice,\mH_\bob
  \cong \C^{V}$, is defined in Definition~\ref{def:EPR} and $\ket{\aux}\in
  \mH_{\alice'} \otimes \mH_{\bob'}$ is arbitrary.
  For each~$x \in \cal{X}$, let $\{M^{x}_a\}_{a \in \cal{A}}$ be a measurement
  on $\mH_{\alice} \otimes \mH_{\alice'}$.
  Let $W\in\{X,Z\}$.
  Then the following are equivalent on the state $\ket{\psi}$:
	\begin{itemize}
	\item On average over $x\sim D$, 
    \begin{equation*}
      \big[ M^x_a \,,\, \big(\tau^W_{[L_x(\cdot)=y]} \otimes \Id_{\overline{V_x}}
      \otimes \Id_{\alice'}\big) \big]  \otimes \Id_\bob \approx_\eps 0\;.
    \end{equation*}
  \item On average over $x \sim D$ and $v$ drawn uniformly from $\ker
    (L_x)^\perp$,
    \begin{equation*}
      \big[ M^x_a \,,\, \big( \tau^W(v) \otimes \Id_{\overline{V_x}} \otimes
      \Id_{\alice'} \big) \big]  \otimes \Id_\bob \approx_\eps 0\;.
    \end{equation*}
	\end{itemize}
	In the equations above, for each $x \in \cal{X}$ we decompose $\mH_\alice = \C^{V_x} \otimes \C^{\overline{V_x}}$ and interpret $\tau^W_{[L_x(\cdot)=y]}$ and $\tau^W(v)$ for $v \in \ker(L_x)^\perp$ as acting on $\C^{V_x}$. 
\end{lemma}

\begin{proof}
  For $x\in \cal{X}$, $a\in \cal{A}$, $y$ in the range of $L_x$, and $v\in \F^k$
  define
  \begin{equation*}
    \Delta^x_{a,y} = \big[ M^x_a \,,\, \big(\tau^W_{[L_x(\cdot)=y]} \otimes
    \Id_{\overline{V_x}} \otimes \Id_{\alice'}\big) \big]\otimes \Id_\bob\:,
  \end{equation*}
  \begin{equation*}
    \Delta^x_a(v) = \big[ M^x_a \,,\, \big( \tau^W(v) \otimes
    \Id_{\overline{V_x}} \otimes \Id_{\alice'} \big) \big]  \otimes \Id_\bob\:.
  \end{equation*}
  By the first item of Lemma~\ref{lem:why-didnt-i-think-of-this-before}, for
  each $v \in \ker(L_x)^\perp$,
  \begin{equation*}
    \Delta^x_a(v) = \sum_{y \in V_x} \omega^{\tr(u_y \cdot v)} \Delta^x_{a,y}\;,
  \end{equation*}
	where for every $y$ in the range of $L_x$, $u_y$ is a fixed element in
  $L_x^{-1}(y)$.
  The expression for the closeness of $\Delta^x_a$ to $0$ on average over $x
  \sim D$ and $v \sim \ker(L_x)^\perp$ (i.e.\ the second quantity of the Lemma
  statement), is equal to
  \begin{equation*}
    \E_{x\sim D} \E_{v \sim \ker(L_x)^\perp} \sum_a\, \bigl \| \Delta^x_a (v)
  \ket{\psi}\bigr \|^2,
  \end{equation*}
  and can be expanded as
  \begin{equation}\label{eq:lemW-1}
    \E_{x\sim D} \E_{v \sim \ker(L_x)^\perp} \sum_a\, \bra{\psi} \sum_{y,y'}
    \omega^{\tr ((u_y'-u_y) \cdot v)} \bigl( \Delta^x_{a,y} \bigr)^\dagger
    \Delta^x_{a,y'} \ket{\psi}\;.
  \end{equation}
  If $y \neq y'$, then by definition $L_x(u_y) \neq L_x(u_{y'})$, and so $u_y' -
  u_y$ is not in $\ker(L_x)$.
  Lemma~\ref{lem:cancellation} (with $V = \ker(L_x)^\perp$) thus implies
  that~\eqref{eq:lemW-1} equals
  \begin{equation*}
    \E_{x\sim D} \sum_a\, \bra{\psi} \sum_{y} \bigl( \bigl( \Delta^x_{a,y}
    \bigr)^\dagger \Delta^x_{a,y} \bigr) \ket{\psi} = \E_{x\sim D} \sum_{a, y}
    \, \bigl \|\Delta^x_{a,y} \ket{\psi} \bigr \|^2
  \end{equation*}
  which is the closeness of $\Delta^x_{a,y}$ to $0$ on average over $x \sim D$.
  Thus, $\Delta^x_a(v) \approx_\eps 0$ on average over $x$ and $v$ if and only
  if $\Delta^x_{a,y} \approx_\eps 0$ on average over $x$.
\end{proof}

\begin{lemma}[Commuting implies twirl]
  \label{lem:mixing-U}
	Let $\cal{X}$ be a finite set and $D$ a distribution on $\cal{X}$. 
  For each~$x\in \cal{X}$, let $\{ M^x_a \}$ be a POVM on $\mH_\alice$, and let
  $\mu_x$ be a distribution over unitary matrices acting on $\mH_\alice$.
  Suppose that on average over $x \sim D$ and $U \sim \mu_x$ it holds that
  \begin{equation*}
    \bigl[ M^x_a , U^\dagger \bigr] \otimes \Id_\bob \approx_{\eps} 0\;,
  \end{equation*}
  where the commutator is evaluated on a state $\ket{\psi} \in \mH_\alice
  \otimes \mH_\bob$.
  Then on average over $x\sim D$,
  \begin{equation*}
    M^x_a \otimes \Id_\bob \approx_\eps \mathscr{T}_{\mu_x} ( M^x_a ) \otimes
    \Id_\bob \;.
  \end{equation*}
\end{lemma}

\begin{proof}
  Observe that
  \begin{equation}
    \E_{x\sim D} \sum_a \, \bigl\| \bigl( \mathscr{T}_{\mu_x} (M^x_a) - M^x_a
    \bigr) \otimes \Id_\bob \ket{\psi} \bigr\|^2 =
    \E_{x\sim D} \sum_a \, \bigl\| \E_{U\sim \mu_x} \bigl( U \big[ M^x_a, U^\dagger
    \bigr] \bigr) \otimes \Id_\bob \ket{\psi} \bigr\|^2.
    \label{eq:twirl2}
  \end{equation}
	Applying Jensen's inequality, the right-hand side of~\eqref{eq:twirl2} is at
  most
  \begin{equation*}
    \E_{x\sim D} \E_{U\sim \mu_x} \sum_a \, \bigl\| \bigl( U \big[ M^x_a, U^\dagger
    \bigr] \bigr) \otimes \Id_\bob \ket{\psi} \bigr\|^2
    = \E_{x\sim D} \E_{U\sim \mu_x} \sum_a \, \bigl\| \big[ M^x_a, U^\dagger
    \bigr]  \otimes \Id_\bob \ket{\psi} \bigr\|^2,
  \end{equation*}
  using the unitary invariance of the Euclidean norm.
  This last quantity is $O(\eps)$, by assumption.
\end{proof}

\begin{lemma}[Commuting with each implies commuting with both]
  \label{lem:commute-with-both}
  Let $\cal{X}$ be a finite set and $D$ a distribution on $\cal{X}$.
  For each~$x\in \cal{X}$, let $\{ M^x_a \}$ be a POVM on $\mH_\alice$ and let
  $\mu_{x,1}, \mu_{x,2}$ be two distributions over unitary matrices acting on
  $\mH_\alice$.
  Suppose that for each $i \in \{1, 2\}$, on average over $x \sim D$ and $U_i
  \sim \mu_{x,i}$,
  \begin{equation}\label{eq:to-be-tagged}
    \bigl[ M^x_a , U_i^\dagger \bigr] \otimes \Id_\bob \approx_{\eps} 0,
  \end{equation}
  where the expression is evaluated on some state $\ket{\psi} \in
  \mH_\alice\otimes\mH_\bob$, where $\mH_\bob \cong \mH_\alice$.
  Suppose further that on average over $x \sim D$ and $U_2 \sim \mu_{x, 2}$,
  \begin{equation}\label{eq:to-be-tagged-dos}
    U_2^\dagger \otimes \Id_\bob \approx_{\eps} \Id_A \otimes U_2\;.
  \end{equation}
  (The corresponding statement for $i = 1$ is not needed.)
  Then on average over $x \sim D$, $U_1 \sim \mu_{x, 1}$, and $U_2 \sim \mu_{x,
    2}$,
  \begin{equation*}
    \bigl[ M^x_a , U_1^\dagger U_2^\dagger \bigr] \otimes \Id_\bob
    \approx_{\eps} 0\;.
  \end{equation*} 
\end{lemma}

\begin{proof}
  The claim follows from the following sequence of approximations:
  \begin{align*}
    M^x_a U_1^\dagger U_2^\dagger \otimes \Id_\bob
    & \approx_{\eps} M^x_a U_1^\dagger \otimes U_2
      \tag{by~\eqref{eq:to-be-tagged-dos}}\\
    & \approx_{\eps} U_1^\dagger M^x_a  \otimes U_2
      \tag{by~\eqref{eq:to-be-tagged} for $i = 1$}\\
    & \approx_{\eps} U_1^\dagger M^x_a U_2^\dagger \otimes \Id_\bob
      \tag{by~\eqref{eq:to-be-tagged-dos}}\\
    & \approx_{\eps} U_1^\dagger U_2^\dagger M^x_a \otimes \Id_\bob
      \tag{by~\eqref{eq:to-be-tagged} for $i = 2$}\;,
  \end{align*}
  where each step also uses Fact~\ref{fact:add-a-proj}.
  This is equivalent to $\bigl[ M^x_a , U_1^\dagger U_2^\dagger \bigr] \otimes
  \Id_\bob \approx_{\eps} 0$, completing the proof.
\end{proof}

The following is a slight generalization of~\cite[Fact~4.25]{NW19},
and we give a similar proof.

\begin{lemma}[Close to sub-measurement implies close to measurement]
  \label{lem:replace-with-measurement}
  Let $\cal{X},\cal{A}$ be finite sets and $D$ be a distribution on $\cal{X}$.
  Suppose that for each $x\in \cal{X}$, $\{A^x_a\}_{a\in \cal{A}}$ is a
  projective measurement and $\{B^x_a\}_{a\in \cal{A}}$ is a set of matrices
  such that each~$B^x_a$ is positive semidefinite and $\sum_{a}B^x_a \leq \Id$.
  Suppose $\{C^x_a\}_{a \in \cal{A}}$ is a POVM such that $C^x_a \geq B^x_a$ for
  all~$x$ and~$a$.
  Then if, on average over $x\sim D$, $A^x_a \approx_{\eps} B^x_a$, then, on
  average over $x\sim D$, $A^x_a \approx_{\eps^{1/2}} C^x_a$.
\end{lemma}

\begin{proof}
  By the triangle inequality,
  \begin{equation*}
    \E_{x} \sum_a \Vert (A^{x}_a - C^{x}_a) \ket{\psi} \Vert^2
    \leq 2\E_{x} \sum_a \Vert (A^{x}_a - B^{x}_a) \ket{\psi} \Vert^2
    + 2 \E_{x}\ \sum_a \Vert (C^x_a - B^x_a) \ket{\psi}\Vert^2.
\end{equation*}
The first term on the right-hand side is $O(\eps)$ by assumption.
For the second,
\begin{multline*}
  \E_{x}\ \sum_a \Vert (C^x_a - B^x_a) \ket{\psi}\Vert^2
  = \E_{x}\ \sum_a \bra{\psi} (C^x_a - B^x_a)^2 \ket{\psi}
  \leq \E_{x}\ \sum_a \bra{\psi} (C^x_a - B^x_a) \ket{\psi}\\
  = 1- \E_{x} \sum_a\bra{\psi}  B^{x}_a \ket{\psi}
  \leq1- \E_{x} \sum_a\bra{\psi}  (B^{x}_a)^2 \ket{\psi}\;,
\end{multline*}
where the middle inequality uses $0\leq C^x_a-B^x_a\leq \Id$ for all $x,a$. 
Write $1 = \E_{x} \sum_a \bra{\psi} (A^{x}_a)^2 \ket{\psi}$, which holds
because~$A$ is a projective measurement.
Then
\begin{multline*}
  \E_{x} \sum_a \bra{\psi} ((A^{x}_a)^2  - (B^{x}_a)^2)\ket{\psi}
  = \Re\Big(\E_{x} \sum_a \bra{\psi} (A^{x}_a  + B^{x}_a)(A^{x}_a - B^{x}_a)
  \ket{\psi}\Big)\\
  \leq \E_{x} \sqrt{\sum_a \Vert (A^{x}_a + B^{x}_a) \ket{\psi}\Vert^2}
	\cdot \sqrt{\sum_a \Vert (A^{x}_a - B^{x}_a)\ket{\psi} \Vert^2}
\end{multline*}
where the first equality follows from the fact that $A^x_a$ and $B^x_a$ are
Hermitian.
For each~$x\in X$ the first square root is $O(1)$.
This allows us to move the expectation into the second square root by Jensen's
inequality.
The result is $O(\eps^{1/2})$ by assumption. 
\end{proof}

Now we put everything together to show the main result of this section. 

\begin{lemma}
  \label{lem:mixing}
  Let~$\cal{X},\cal{A}$ be a finite sets and~$D$ be a distribution over
  $\cal{X}$.
  For each~$x \in \cal{X}$, let $V_x$ be a register subspace of $V = \F^k$, let
  $U_x$ be a register subspace of~$V_x$, and let $L_x:U_x \to U_x$ be a linear
  map.

  Consider a state $\ket{\psi} = \ket{\EPR}_V \otimes \ket{\aux}$, where
  $\ket{\EPR}_{V} \in \mH_\alice \otimes \mH_\bob$, for $\mH_\alice,\mH_\bob
  \cong \C^{V}$, is defined in Definition~\ref{def:EPR} and $\ket{\aux}\in
  \mH_{\alice'} \otimes \mH_{\bob'}$ is arbitrary.
  For each~$x \in \cal{X}$, let $\{M^{x}_{y,a}\}_{y \in U_x,\, a\in \cal{A}}$ be
  a projective measurement on $\mH_{V_x} \otimes \mH_{\alice'}$.
  Suppose that on average over~$x \sim D$ the following conditions hold.
  \begin{alignat*}{2}
    \text{(Consistency):} && \quad
    \bigl( M^{x}_{y} \otimes \Id_{\overline{V_x}} - \tau^Z_{[L_x(\cdot) = y]}
    \otimes \Id_{\overline{U_x}} \otimes \Id_{\alice'}\bigr) \otimes \Id_\bob
    \approx_\eps 0\;, &\\
    \text{(Commutation):} && \quad
      \bigl[ M^{x}_{y,\, a} \otimes \Id_{\overline{V_x}}\,,\, \tau^Z_{z} \otimes
      \Id_{\overline{U_x}}   \otimes \Id_{\alice'} \bigr] \otimes
      \Id_\bob\approx_\eps 0\;, &\\
      && \quad
      \bigl[ M^{x}_{y,\, a}\otimes \Id_{\overline{V_x}}\,,\,
      \tau^X_{[L_x^\perp(\cdot) = y^\perp]}  \otimes \Id_{\overline{U_x}}
      \otimes \Id_{\alice'} \bigr]  \otimes \Id_\bob  \approx_\eps 0\;.&
  \end{alignat*}
  Here, the projector $\tau^Z_{[L_x(\cdot) = y]}$ acts on the register subspace
  $\mH_{U_x}$, and $\overline{U_x}$ and $\overline{V_x}$ denote the
  complementary register subspaces of $U_x$ and $V_x$, respectively, within $V$. Furthermore, the answer summations are over all $y, z \in U_x$ and $a \in \cal{A}$.
  
  
  Then for each $x\in \cal{X}$ and $y \in U_x$, there exists a POVM measurement
  $\bigl\{ M^{x,\, y}_a \bigr\}_{a \in \cal{A}} $ on $\mH_{V_x \setminus
    U_x}\otimes \mH_{\alice'}$ such that on average over $x \sim D$,
  \begin{equation*}
    (M^{x}_{y,\, a} \otimes \Id_{\overline{V_x}}) \otimes \Id_{B}
    \approx_{\eps^{1/2}} \Bigl(\tau^Z_{[L_x(\cdot) = y]} \otimes M^{x,\, y}_a
    \otimes \Id_{\overline{V_x}}\Bigr) \otimes \Id_\bob\;.
  \end{equation*}
\end{lemma}

\begin{proof}
  For each~$x \in X$, let $\zeta_{x}$ be the uniform distribution over $U_x$ and
  $\chi_{x}$ be the uniform distribution over $\ker(L_x)$.
  \hnote{In response to Bowen comments from second file, item 15c, elaborating on how Lemma 8.12 is applied.} Applying \Cref{lem:W} to the first commutation assumption (letting ``$L_x$'' in \Cref{lem:W} be the identity map on $U_x$, and letting ``$D$'' be the uniform distribution over $\ker(L_x)^\perp = U_x$), we get that on average over $x \sim \cal{X}$ and $v \sim \zeta_x$, we have
  \begin{equation*}
    \bigl[ M^{x}_{y,\, a} \otimes \Id_{\overline{V_x}} \,,\, \tau^Z(v) \otimes
    \Id_{\overline{U_x}} \otimes \Id_{\alice'} \bigr]  \otimes \Id_\bob
    \approx_\eps 0\;.
  \end{equation*}
  Applying \Cref{lem:W} to the second commutation assumption (letting ``$L_x$'' in \Cref{lem:W} be the $L_x^\perp$ in the statement of \Cref{lem:mixing} and letting $D$ be the uniform distribution over $\ker(L_x^\perp)^\perp$, which is $\ker(L_x)$ by Lemma~\ref{lem:L_perp_perp}), we get 
on average over $x \sim \cal{X}$ and $v \sim \xi_x$, we have
  \begin{equation*}
    \bigl[ M^{x}_{y,\, a} \otimes \Id_{\overline{V_x}} \,,\, \tau^X(v) \otimes
    \Id_{\overline{U_x}} \otimes \Id_{\alice'} \bigr]  \otimes \Id_\bob
    \approx_\eps 0\;.
  \end{equation*}    
By Lemma~\ref{lem:commute-with-both}, this implies that on average over $x
  \sim D$, $u \sim \zeta_{x}$, and~$v \sim \chi_{x}$,
  \begin{equation*}
    \bigl[ M^{x}_{y,\, a} \otimes \Id_{\overline{V_x}} \,,\, \tau^Z(u) \tau^X(v)
    \otimes \Id_{\overline{U_x}} \otimes \Id_{\alice'} \bigr] \otimes \Id_\bob
    \approx_\eps 0\;.
  \end{equation*}
  By Lemma~\ref{lem:mixing-U} and Lemma~\ref{lem:mixing-L},
  this implies that on average over $x\sim D$,
  \begin{align}
    M^{x}_{y,\, a} \otimes \Id_{\overline{V_x}} \otimes \Id_\bob
    & \approx_\eps \bigl( \mathscr{T}_{\chi_{x}} \circ \mathscr{T}_{\zeta_{x}} (
      M^{x}_{y,\, a}) \bigr) \otimes \Id_{\overline{V_x}} \otimes
      \Id_\bob\notag\\
    & = \Big(\sum_{y' \in V_x} \tau^Z_{[L_x(\cdot) = y']} \otimes  M^{x,\,
      y'}_{y,\, a} \Big) \otimes \Id_{\overline{V_x}} \otimes
      \Id_\bob\;,\label{eq:mixing-eq-1}
  \end{align}
	for some POVM measurement $\{M^{x,\, y'}_{y,\, a} \}$ on $\mH_{V_x \setminus
    U_x} \otimes \mH_{\alice'}$.
	
	In the following sequence of equations, whenever an operator does not act on a
  subsystem it should be assumed that it is appropriately tensored with the
  identity.
  For clarity, we explicitly indicate using a subscript $\alice$ or $\bob$
  whether a Pauli operator acts on $\mH_\alice$ or $\mH_\bob$.
  Then on average over $x\sim D$ we have
  \begin{align*}
    M^{x}_{y, \, a}
    & =  M^{x}_{y, \, a} \cdot M^{x}_y  \tag{$M^x$ is projective}\\
    & \approx_\eps  M^{x}_{y, \, a} \cdot \big(
      \tau^Z_{[L_x(\cdot)=y]}\big)_\alice
      \tag{Consistency assumption}\\
    & \approx_0 M^x_{y, \, a} \otimes \big(\tau^Z_{[L_x(\cdot)=y]}\big)_\bob
      \tag{Paulis are self-consistent}\\
    & \approx_\eps \Big( \sum_{y'} \big(\tau^Z_{[L_x(\cdot) = y']}\big)_\alice
      \otimes  M^{x,\, y'}_{y,\, a} \Big) \otimes \big(\tau^Z_{[L_x(\cdot)=y]}
      \big)_\bob \tag{Equation~\eqref{eq:mixing-eq-1}}\\
    & \approx_0 \sum_{y'} \big(\tau^Z_{[L_x(\cdot) =
      y']}\tau^Z_{[L_x(\cdot)=y]}\big)_\alice \otimes  M^{x,\, y'}_{y,\, a}
      \tag{Paulis are self-consistent} \\
    & = \big( \tau^Z_{[L_x(\cdot) = y]} \big)_\alice
      \otimes M^{x,\, y}_{y,\, a} \;,
  \end{align*}
  where $\approx_0$ indicates equality with respect to the state $\ket{\EPR}_{V}
  \otimes \ket{\aux}$, and we have repeatedly used Fact~\ref{fact:add-a-proj}.
  This is essentially the statement promised by the lemma, except that
  $\{M^{x,\, y}_{y,\, a}\}_a$ does not necessarily sum to identity (since we
  only sum over~$a$).
  To remedy this, define $M^{x,\, y}_a = \sum_{y'} M^{x,\, y}_{y',\, a}$ and note
  that $M^{x,\, y}_{a} \geq M^{x,\, y}_{y,\, a}$, so
  by~Lemma~\ref{lem:replace-with-measurement} and the fact that~$M^x$ is
  projective,
  \begin{equation}\label{eq:mixing-eq-2}
    (M^{x}_{y,\, a} \otimes \Id_{\overline{V_x}}) \otimes \Id_{B} \approx_{\eps^{1/2}}
    \Bigl(\tau^Z_{[L_x(\cdot) = y]} \otimes M^{x,\, y}_a \otimes \Id_{\overline{V_x}}\Bigr)
    \otimes \Id_\bob\;.
  \end{equation}
  $\{M^{x\,y}_a\}$ is the POVM measurement guaranteed in the lemma statement,
  which concludes the proof.
\end{proof}


\subsubsection{Preliminary lemmas}

We show a few simple lemmas that allow us to argue about measurements that
have a decomposition across a tensor product of two Hilbert spaces, within the
space of a single player.

\def\ha{\mH_{\alice}}
\def\hb{\mH_{\bob}}
\def\haq{\mH_{\alice,\, a}}
\def\haa{\mH_{\alice',\, a}}
\def\hbq{\mH_{\bob,\, a}}
\def\hba{\mH_{\bob',\, a}}
\def\sq{\ket{\psi_{\textsc{que},\, a}}}
\def\sa{\ket{\psi_{\textsc{ans},\, a}}}

\begin{lemma}
  \label{lem:conditional}
	Let $\cal{A}, \cal{B}$ be finite sets.
	Let $\ket{\psi} \in \ha \otimes \hb$ be a state.
  Consider the following: for all $a \in \cal{A}$,
  \begin{enumerate}
  \item Let $\haq$, $\hbq$, $\haa$, and $\hba$ be Hilbert spaces such that
    \begin{equation*}
  		\ha = \haq \otimes \haa\quad
      \text{and}\quad
      \hb = \hbq \otimes \hba\;,
		\end{equation*}
    and let $\sq \in \haq \otimes \hbq$ be a ``question state'' and $\sa \in
    \haa \otimes \hba$ be an ``answer state'' such that $\ket{\psi} = \sq
    \otimes \sa$.
  \item Let $Q_a$ be projectors on $\haq$ such that $\{Q_a \otimes
    \Id_{\haa}\}$ forms a projective measurement on $\ha$ and let
    $\{A^{a}_b\}_{b \in \cal{B}}$ and $\{B^{a}_b\}_{b \in \cal{B}}$ be matrices
    acting on $\haa$.
  \item Let $D$ be the distribution on~$\cal{A}$ obtained by measuring
    $\ket{\psi}$ using $\{ Q_a \otimes I_{\haa}\}_{a \in \cal{A}}$.
  \end{enumerate}
  Then the following are equivalent:
	\begin{itemize}
	\item On average over $a\sim D$ and with respect to state $\ket{\psi}$,
    \begin{equation*}
      (\Id_{\haq} \otimes A^{a}_b) \otimes \Id_\bob \approx_\eps
      (\Id_{\hbq} \otimes B^{a}_b) \otimes \Id_\bob\;.
    \end{equation*}
  \item $\bigl( Q_a \otimes A^{a}_b \bigr) \otimes \Id_{\bob} \approx_\eps
    \bigl( Q_a \otimes B^{a}_b \bigr) \otimes \Id_{\bob}$ on state $\ket{\psi}$.
	\end{itemize}
\end{lemma}

\begin{proof}
 Expand
  \begin{align*}
    & \sum_{a,b} \norm{ (Q_a \otimes A^a_b - Q_a \otimes B^a_b) \otimes
      \Id_{\bob} \cdot \sq \otimes \sa }^2\\
    = & \sum_{a,b} \norm{ Q_a \otimes (A^a_b -  B^a_b) \otimes
        \Id_{\bob} \cdot \sq \otimes \sa}^2\\
    = & \sum_{a, b} \norm{ Q_a \otimes I_{\hbq} \sq}^2 \cdot
        \norm{  (A^a_b -  B^a_b) \otimes \Id_{\hba} \sa }^2\\
	= & \sum_{a, b} \norm{ Q_a \otimes \Id_{\bob} \ket{\psi}}^2 \cdot
      \norm{  (A^a_b -  B^a_b) \otimes \Id_{\hba} \sa }^2\\
    = & \E_{a \sim D} \sum_{b}
        \norm{  (A^a_b -  B^a_b) \otimes \Id_{\hba} \sa }^2\\
    = & \E_{a \sim D} \sum_{b}
      \norm{  \Id_{\haq} \otimes (A^a_b -  B^a_b) \otimes
        \Id_{\bob} \cdot \sq \otimes \sa }^2\;.
  \end{align*}
  In going from the third to the fourth line we used the fact that
  \[
    \norm{ Q_a \otimes I_{\hbq} \sq }^2 =
    \norm{ Q_a \otimes \Id_{\haa} \otimes I_{\bob} \sq \otimes \sa }^2
    = \norm{ Q_a \otimes \Id_{\haa} \otimes \Id_{\bob} \ket{\psi} }^2\;.
  \]
  Hence, the first line is $O(\eps)$ if and only if the last one is.
\end{proof}

\begin{lemma}\label{lem:commutation-simplification}
  Let $\mA,\mB,\mC$ be finite sets.
  Let $\ket{\psi} \in \ha \otimes \hb$ be a state, and let $\{A_{a,\, b}\}_{a
    \in \mA, b \in \cal{B}}$ and $\{B_{a,\, c}\}_{a \in \mA, c \in \cal{C}}$ be
  POVMs acting on $\ha$.
  Suppose further that:
  \begin{enumerate}
  \item The measurements approximately commute,
    i.e.\, \label{item:almost-compute}
    \begin{equation*}
      \bigl[ A_{a,\, b}, B_{a,\, c} \bigr] \otimes \Id_\bob \approx_\delta 0\;,
    \end{equation*}
    where the approximation holds with respect to the state $\ket{\psi}$. 
  \item For all $a \in \cal{A}$, there exist Hilbert spaces $\haq$, $\haa$,
    $\hbq$, $\hba$, and states $\sq \in \haq \otimes
    \hbq$, $\sa \in \haa \otimes \hba$ such that
    \begin{align*}
      \ha & = \haq \otimes \haa\;,\\
      \hb & = \hbq \otimes \hba\;,\\
      \ket{\psi} & = \sq \otimes \sa\;.
    \end{align*}
  \item \label{item:weird-decomposition} For all $a \in \cal{A}$, there exist
    projectors $Q_a$ on $\haq$ and matrices $\{A^{a}_b\}_{b \in
      \cal{B}}$, $\{B^{a}_c\}_{c \in \cal{C}}$ acting on $\haa$
    such that $\{Q_a \otimes \Id_{\haa}\}$ is a projective
    measurement on $\ha$ and
    \begin{equation*}
      A_{a, \, b} = Q_a \otimes A^a_b\;, \quad
      B_{a, \, c} = Q_a \otimes B^a_c\;.
    \end{equation*}
  \end{enumerate}
  Then
  \begin{equation*}
    \bigl[ \Id_{\haq} \otimes A^{a}_b \,,\, \Id_{\haq}
    \otimes B^{a}_c \bigr] \otimes \Id_\bob \approx_\delta 0\;,
  \end{equation*}
  on average over $a \sim D$ where $D$ is the distribution on~$\cal{A}$ obtained
  by measuring $\ket{\psi}$ using $\{ Q_a \otimes I_{\haa} \}_{a \in \cal{A}}$.
\end{lemma}
 
\begin{proof}
  The assumptions of the lemma imply that
  \begin{align*}
    (Q_a \otimes A^a_b B^a_c) \otimes \Id_\bob
    & = ((Q_a \otimes A^a_b) \cdot (Q_a \otimes B^a_c)) \otimes \Id_\bob
      \tag{$Q_a$ is a projector}\\
    & = (A_{a, b} \cdot B_{a, c}) \otimes \Id_\bob
      \tag{Item~\ref{item:weird-decomposition}}\\
    & \approx_\delta (B_{a, c} \cdot A_{a, b}) \otimes \Id_\bob
      \tag{Item~\ref{item:almost-compute}}\\
    & = ((Q_a \otimes B^a_c) \cdot (Q_a \otimes A^a_b)) \otimes \Id_\bob
      \tag{Item~\ref{item:weird-decomposition}}\\
    & = (Q_a \otimes B^a_c A^a_b) \otimes \Id_\bob\;.
      \tag{$Q_a$ is a projector}
  \end{align*}
  We apply Lemma~\ref{lem:conditional} as follows.
  The set ``$\mA$'' in Lemma~\ref{lem:conditional} is the same as $\cal{A}$
  here, and the set ``$\mB$'' is the product set $\mB \times \mC$ here.
  The matrices ``$\{A^a_b\}$'' are $\{A^a_{b}B^a_c\}$ here and ``$\{B^{a}_b\}$''
  are $\{B^a_c A^a_{b}\}$ here.
  We then obtain, on average over $a \sim
  D$,
  \begin{equation*}
    (\Id_{\mH_{\alice, a}}\otimes A^{a}_b B^a_c) \otimes \Id_\bob
    \approx_\delta (\Id_{\mH_{\alice, a}} \otimes B^{a}_c A^a_b) \otimes \Id_\bob\;.
  \end{equation*}
  This implies the conclusion of the lemma.
\end{proof}

\begin{lemma}
  \label{lem:conditional-consistency}
	Let $ \cal{Y}$ be a finite set and for all $y\in \cal{Y}$ let $\{A^y_{x,\, z}\}$ be a POVM on $\mH_\alice$. Let 
  $\{B_{x,\, y,\, z}\}$ be a projective measurement on $\mH_\bob$.
  Suppose that
  \begin{equation}\label{eq:cc-1}
    \sum_{x,\, y,\, z} \bra{\psi} A^y_{x,\, z} \otimes B_{x,\, y,\, z}
    \ket{\psi} \ge 1 - \delta\;.
  \end{equation}
  Then with respect to state $\ket{\psi}$,
  \begin{equation*}
    \ia \otimes B_{x,\, y,\, z} \approx_\delta A^y_{x,\, z} \otimes B_{x,\,
      y}\;.
  \end{equation*}
\end{lemma}

\begin{proof}
  Using the fact that $\{B_{x,\, y,\, z}\}$ is projective, we have $B_{x,\, y,\, z} =
  B_{x,\, y} B_z$ for all $x,y,z$, so that~\eqref{eq:cc-1} implies
  \begin{equation*}
    \sum_{x,\, y,\, z} \bra{\psi} A^y_{x,\, z} \otimes B_{x,\, y} B_{z}
    \ket{\psi} \ge 1 - \delta\;.
  \end{equation*}
  Define $\hat{A}_z = \sum_{x,\, y} A^y_{x,\, z} \otimes B_{x,\, y}$ and
  $\hat{B}_z = \ia \otimes B_z$.
  The above equation simplifies to
  \begin{equation*}
    \sum_{y,\, z} \bra{\psi} \hat{A}_z \hat{B}_z \ket{\psi} \ge 1 - \delta\;.
  \end{equation*}
  This implies that $\hat{A}_z \approx_\delta \hat{B}_z$ as
  \begin{equation*}
    \sum_z \norm{(\hat{A}_z - \hat{B}_z) \ket{\psi}}^2 = \sum_z \bra{\psi}
    \hat{A}_z^2 + \hat{B}_z^2 \ket{\psi} - 2 \bra{\psi} \hat{A}_z \hat{B}_z
    \ket{\psi} \le 2 \delta\;,
  \end{equation*}
  where the equality uses the fact that $\hat{A}_z$ and $\hat{B}_z$ commute.
  To conclude the proof, we have
  \begin{equation*}
    \ia \otimes B_{x,\, y,\, z} = \ia \otimes B_{x,\, y} B_z \approx_\delta (\ia
    \otimes B_{x,\, y}) \sum_{x',\, y'} A^{y'}_{x',\, z} \otimes B_{x',\, y'} =
    A^y_{x,\, z} \otimes B_{x,\, y},
  \end{equation*}
  where the approximation follows from $\hat{A}_z \approx_\delta \hat{B}_z$ and
  \cref{fact:add-a-proj}.

\end{proof}

\subsubsection{Soundness proof}

We analyze the soundness of the introspective verifier.

\begin{proof}[Proof of the soundness part of Theorem~\ref{thm:introspection}]
  Let $\verifier=(\sampler, \decider)$ be a normal form verifier such that
  $\sampler$ is an $\ell$-level sampler.
  Recall the definition of the introspective verifier $\tvint =
  (\tsint, \tdint)$ corresponding to $\verifier$ from
  Section~\ref{sec:intro-verifier}.
  Fix an index $n\geq 1$ and let $N = 2^n$.
  Let $R = N^\lambda$ and $\introparams(R) = (q,m,d)$,
  as in Section~\ref{sec:intro-verifier}.
 
  As in the proof of completeness part of \cref{thm:introspection}, we make the
  following simplifications.
  First, we analyze the soundness of the typed introspective verifier
  $\tvint$; the soundness of the detyped verifier
  $\vint = \detype(\tvint)$ follows from Lemma~\ref{lem:detyping-verifiers}
  and the fact that the type set $\type^\intro$ has size $O(\ell)$.
  Second, analogously to Remark~\ref{rk:downsizing} without loss of generality
  for notational simplicity we consider strategies for questions sampled from
  $\tilde{\sampler}^\intro$ rather than the downsized sampler
  ${\sampler}^\intro$.

  Suppose that $\val^*(\verifier_n^\intro) > 1-\eps$ for some $0 < \eps < 1$,
  and let $\strategy = (\ket{\hat{\psi}}, \hat{A}, \hat{B})$ be a strategy for
  $\tvint_n$ with value at least~$1-\eps$.
  Since $\strategy$ has success probability that is strictly positive, the
  decider $\tdint$ does not automatically reject, which means that
	\begin{equation}\label{eq:intro-sound-lambda}
	s(N) \,\leq\, N^\lambda\;.
	\end{equation}
	We analyze each of the tests performed by $\tdint$ (see
  Figure~\ref{fig:intro-decider}) in sequence, and state consequences of each
  test.
  We start with Item~\ref{enu:pauli}, the Pauli test.
	
  \begin{lemma}\label{lem:intro-pauli-strat}
    There exists a function $\delta_1(\eps,R) = a( (\log R)^a \eps^{b} + (\log R)^{-b})$ for universal constants $a > 1$ and $0 < b < 1$ such that the following holds.
    Let $Q = M \log q$, where $M=2^m$.
    There is a projective strategy $\strategy'=(\ket{\psi},A,B)$ for
    $\tvint_n$ that succeeds with probability at least $1-\delta_1$
    and furthermore
    \begin{equation}
      \label{eq:intro-state-a}
      \ket{\psi}_{\alice\bob} =  \ket{\EPR_2}_{\alice'\bob'}^{\otimes Q} \otimes
      \ket{\aux}_{\alice''\bob''}
    \end{equation}
    for some  bipartite state $\ket{\aux}$, and for all $W\in\{X,Z\}$,
    \begin{equation}\label{eq:intro-sound-paulia}
      A^{\Pauli,W}_x = \sigma^W_x \;,\qquad B^{\Pauli,W}_x = \sigma^W_x \;,
    \end{equation}
    where $\sigma^W_x$ acts on the first $s(N)$ qubits of player $\alice$'s
    share (resp.
    $\bob$'s share) of $\ket{\EPR_2}^{\otimes Q}$.
  \end{lemma}
	
  \begin{proof}
    Given the definition of the type graph $G^\intro$, for
    $((\hat{t}_\alice,\hat{x}_\alice),(\hat{t}_\bob,\hat{x}_\bob))$ sampled
    according to $\mu_{\tsint,n}$ it holds that
    $(\hat{t}_\alice,\hat{t}_\bob) \in \type^\pauli \times \type^\pauli$ with
    constant probability.
    Therefore, conditioned on the Pauli test, Item~\ref{enu:pauli}, being
    executed, $\strategy$ must succeed in the test with probability $1-O(\eps)$.

    Observe that conditioned on the test being executed, the distribution of
    $((\hat{t}_\alice,\hat{x}_\alice),(\hat{t}_\bob,\hat{x}_\bob))$ is, by
    definition, exactly the distribution of questions in the Pauli basis game
    with parameters $\qldparams$, as described in Section~\ref{sec:qld-game}.
    By \Cref{cor:pauli-binary} it follows that there exists a local isometry
    $\phi = \phi_\alice \otimes \phi_\bob$ and a state $\ket{\aux}\in
    \mH_{\alice''}\otimes \mH_{\bob''}$ such that
    \begin{equation}\label{eq:intro-sound-0}
      \Vert \phi (\ket{\hat{\psi}}) - \ket{\EPR_2}^{\otimes Q} \otimes
      \ket{\aux} \Vert^2 \leq \delta'(\eps,R)\;,
    \end{equation}
    where $\delta'(\eps,R)$ is an upper bound on $\delta_\qld(O(\eps),q,m,d)$
    that only depends on $\eps$ and $R$, as stated in
    Lemma~\ref{lem:delta-bound}.
    In addition, defining $A^{\hat{x}}_{\hat{a}} = \phi_\alice~
    \hat{A}^{\hat{x}}_{\hat{a}}~\phi_\alice^\dagger$ for all questions $\hat{x}$ and answers
    $\hat{a}$, for $W \in \{X, Z\}$ it holds that
    \begin{equation}
      \label{eq:intro-sound-pauli}
      A^{\Pauli,W}_x \otimes \Id_\bob \approx_{\delta'(\eps,R)}
      \sigma^W_x \otimes \Id_\bob\;,
    \end{equation}
    and a similar set of equations hold for operators associated with the second
    player.
    Using Naimark's theorem as formulated in~\cite[Theorem 5.1]{ML20}, at the
    cost of extending the state $\ket{\aux}$ we may assume that the measurements
    are projective without loss of generality.
    Define the strategy $\strategy'$ which uses the state $\ket{\EPR_2}^{\otimes
      Q} \otimes \ket{\aux}$ and measurement operators $\{ A^{\hat{x}}_{\hat{a}}
    \}$ and $\{ B^{\hat{x}}_{\hat{a}} \}$ for all questions $\hat{x}$, except
    for $(\Pauli,W)$-type questions where instead the Pauli measurements
    $\sigma^W_{\hat{a}}$ are used.
    Using~\eqref{eq:intro-sound-0} and~\eqref{eq:intro-sound-pauli} the strategy
    $\strategy'$ succeeds in $\tvint_n$ with probability at least
    $1-\delta'(\eps,R)$.
	
    The claimed bound on $\delta_1$ follows from the bound given in
    Lemma~\ref{lem:delta-bound}.
  \end{proof}
	
  In the remainder of the proof we analyze the strategy $\strategy'$ from
  Lemma~\ref{lem:intro-pauli-strat}.
  We use the following notation conventions:
  \begin{enumerate}
  \item We use indices $\alice$ and $\bob$ to label each player's Hilbert space
    after application of the isometry $\phi$ from
    Lemma~\ref{lem:intro-pauli-strat}.

  \item We write~$V$ for the register subspace of $\F_2^{Q}$ spanned by $e_1,
    \ldots, e_{s(N)}$ and $\overline{V}$ for its complement.
    (Note that by definition of $\introparams$ in
    Section~\ref{sec:qld-complexity} it holds that $s(N)\leq R \leq Q$, where
    the first inequality follows from~\eqref{eq:intro-sound-lambda}.)

  \item Whenever we write a Pauli operator $\sigma^W_a$ the register on which
    the operator acts should always be clear from context, and is implicit from
    the space in which the outcome $a$ ranges.
  
  \item For measurement operators in the introspection game, the variables for
    the measurement outcomes follow the specification of the ``answer key'' in
    Figure~\ref{fig:decider_pauli} (for $\type^\Pauli$-type questions) and
    Figure~\ref{fig:intro-decider} (for all other question types).
    For example, the measurement operators $\{A_x^{\Pauli,W} \}$ corresponding
    to question type $(\Pauli,W)$ are indexed by vectors $x \in \F_q^Q$ where $Q
    = M \log q$ where $M = 2^m$.
    The measurement operators corresponding to question type $(\Intro,\trole)$
    for $\trole \in \AB$ are indexed by pairs $(y,a) \in V \times \{0,1\}^{\leq
      3Q}$.\footnote{Technically the answer $a$ may be a binary string of any
      length; however, if $a$ is too long the decider rejects due to the answer
      length check.
      Thus we assume without loss of generality that the answer $a$ is a binary
      string of length at most $3Q = 3\cdot 2^m \cdot \log q$.}
    We often refer to marginalized measurement operators, e.g., the operator
    $A^{\Intro,\trole}_{y}$ denotes marginalizing $A^{\Intro,\trole}_{y,\, a}$ over
    all $a$.
    In these cases, the part of the answer that is marginalized over will be
    clear from context.
		
  \item We use the notation $\delta$ to denote a function which is polynomial in
    $\delta_1$, although the exact expression may differ from occurrence to
    occurrence.
    The polynomial itself may depend on $\ell$, but we leave this dependence
    implicit; due to the use of inductive steps that involve taking the square
    root of the error $\ell$ times in sequence (e.g.
    Lemma~\ref{lem:intro-sound-induction}) the exponent generally depends on
    $\ell$.
  \end{enumerate}

  The next two lemma derive conditions implied by Items~\ref{enu:sampling}
  and~\ref{enu:hiding} of the checks performed by the decider $\tdint$
  described in Figure~\ref{fig:intro-decider}.
  As these two parts are performed independently for the two possible values of
  $\trole\in\{\alice,\bob\}$, we only discuss the case where $\trole = \alice$.
  For notational simplicity, whenever possible we omit $v$ when referring to the
  measurement operators.
  For example $L$, $\ai$ and $\bs$ are used as shorthand notation
  for $L^\alice$, $A^{\Intro,\, \alice}_{y,\, a}$, and $B^{\Sample,\,
    \alice}_{z,\, a}$ respectively.

  \begin{lemma}[Sampling test, Item~\ref{enu:sampling} of
    Figure~\ref{fig:intro-decider}]
    \label{lem:intro-sound-1}
    For each $k\in\{ 1, 2, \ldots, \ell\}$,
    \begin{align}
      \ia \otimes \bs[z]
      & \simeq_\delta \sigma^Z_z \otimes \ib\;,
        \label{eq:intro-sound-1} \\
      \ai[y_{\leq k},\, a] \otimes \ib
      & \simeq_\delta \ia \otimes \bs[\Llek,\, a]\;,
        \label{eq:intro-sound-3b}
    \end{align}
    where $z$ ranges over $V$ and $y_{\leq k}$ ranges over $L_{\leq k}(V)$.
    Moreover, analogous equations hold with operators acting on the other side
    of the tensor product.
  \end{lemma}
	
  \begin{proof}
    When $((\hat{t}_\alice,\hat{x}_\alice),(\hat{t}_\bob,\hat{x}_\bob))$ is
    sampled according to $\mu_{\tsint,n}$, each check in
    Item~\ref{enu:sampling} of Figure~\ref{fig:intro-decider} is executed with
    probability $\Omega(1/\ell)$ (this is due to the number of types in
    $\type^\intro$ and the structure of the type graph $G^\intro$).
    Therefore, in each of the checks specified by Items~\ref{enu:sampling-pauli}
    and~\ref{enu:sampling-intro}, the strategy $\strategy'$ succeeds with
    probability at least $1 - O(\ell \delta)$, conditioned on the test being
    executed.
    Item~\ref{enu:sampling-pauli} for $w=\alice$ combined
    with~\eqref{eq:intro-sound-paulia} implies~\eqref{eq:intro-sound-1}.
    Item~\ref{enu:sampling-intro} for $w=\alice$, combined with
    Fact~\ref{fact:data-processing}, implies~\eqref{eq:intro-sound-3b}.
    The lemma follows from repeating the same argument with the tensor factors
    interchanged.
  \end{proof}

  \begin{lemma}[Hiding test, Item~\ref{enu:hiding} of
    Figure~\ref{fig:intro-decider}]\label{lem:intro-sound-2}
    For each $k \in \{1, \ldots, \ell+1\}$,
    \begin{equation}\label{eq:intro-sound-4letter-before-a}
      \ai[y_{< k},\, a] \abc \br[y_{< k},\, a]\;,
    \end{equation}
    and if $k\leq \ell$,
    \begin{equation}\label{eq:intro-sound-4a-2}
      \ia \otimes \bhk  \approx_\delta \zLlk \otimes \ib\;.
    \end{equation}
    Furthermore, for all $j,k \in \{1,\ldots,\ell -1 \}$ such that $j \leq k$,
    we have
    \begin{equation}
      \label{eq:hide-chain}
      \ahk[y_{<j},\, y_{\leq j}^\perp] \otimes \ib \approx_\delta
      \ahp[y_{<j},\, y_{\leq j}^\perp] \otimes \ib\;.
    \end{equation}
    Analogous equations
    to~\eqref{eq:intro-sound-4letter-before-a},~\eqref{eq:intro-sound-4a-2},
    and~\eqref{eq:hide-chain} hold with operators acting on the other side of
    the tensor product.
  \end{lemma}
	
  \begin{proof}
    When $((\hat{t}_\alice,\hat{x}_\alice),(\hat{t}_\bob,\hat{x}_\bob))$ is
    sampled according to $\mu_{\tsint,n}$, each check in
    Item~\ref{enu:hiding} of Figure~\ref{fig:intro-decider} is executed with
    probability $\Omega(1/\ell)$.
    Therefore, in each of the checks specified by
    Items~\ref{enu:hiding-intro},~\ref{enu:hiding-read}, and~\ref{enu:hiding-same}
    (conditioned on the right types) the strategy
    $\strategy'$ succeeds with probability at least $1 - O(\ell \delta)$.

    Item~\ref{enu:hiding-intro} for $w=\alice$ combined with
    Fact~\ref{fact:data-processing}
    implies~\eqref{eq:intro-sound-4letter-before-a}.
    Combining \cref{eq:intro-sound-1,eq:intro-sound-3b} with
    \eqref{eq:intro-sound-4letter-before-a} yields
    \begin{equation}
      \label{eq:intro-sound-4a}
      \ia \otimes \br[y_{< k}] \approx_\delta \zLlk \otimes \ib\;.
    \end{equation}
    The fact that $\strategy'$ is projective and succeeds in
    Items~\ref{enu:hiding-read} and~\ref{enu:hiding-same} of
    Figure~\ref{fig:intro-decider} with probability $1 - O(\ell \delta)$,
    along with Item~\ref{item:consistency-implies-approx} of
    \Cref{fact:agreement}, imply~\eqref{eq:intro-sound-4a-2}.
  
    We now establish the ``Furthermore'' part of the lemma statement.
    Let $1 \leq j \leq k \leq \ell-1$.
    Item~\ref{enu:hiding-same}, \Cref{item:consistency-implies-approx} of
    \Cref{fact:agreement}, and \Cref{fact:data-processing} imply that
    \begin{equation}
      \ahk[y_{<j},\, y_{\leq j}^\perp] \otimes \ib \approx_\delta
      \ia \otimes \bhp[y_{<j},\, y_{\leq j}^\perp].
    \end{equation}
    Item~\ref{enu:intro-consistency} and \Cref{fact:data-processing} imply
    \begin{equation}
      \ahp[y_{<j},\, y_{\leq j}^\perp] \otimes \ib \simeq_\delta
      \ia \otimes \bhp[y_{<j},\, y_{\leq j}^\perp]\;.
    \end{equation}
    This proves~\eqref{eq:hide-chain}.

    The lemma follows from repeating the same arguments with the tensor factors
    interchanged.
	\end{proof}
	
	We exploit the tests performed in Item~\ref{enu:hiding} further to show the
  following lemma.

  \begin{lemma}
    \label{lem:hide-rigidity}
    For all $k \in \{1, 2, \ldots, \ell\}$,
    \begin{equation*}
      \ia \otimes \bhk[y_{<k},\, y^\perp_k,\, x_{>k}] \approx_\delta
      \bigl( \zLlk \otimes \xLperpk \otimes \x[x_{>k}] \bigr) \otimes \ib\;,
    \end{equation*}
    and an analogous equation holds with operators acting on the other side of
    the tensor product.
  \end{lemma}

  \begin{proof}
    The proof is by induction on $k$. We first show the case $k=1$.
    Under the distribution $\mu_{\tsint,n}$, the check in
    \cref{enu:hiding-pauli} of \cref{fig:intro-decider} is executed with
    probablity $\Omega(1/\ell)$. That part of the check for $w=\alice$ together
    with \cref{eq:intro-sound-paulia} implies that
    \begin{equation}
      \label{eq:base-case-hide}
      \ia \otimes B^{\Hide{1}}_{y_1^\perp,\, x_{>1}} \approx_\delta \bigl(
      \xLperpk[1] \otimes \x[x_{>1}] \bigr) \otimes \ib\;.
    \end{equation}
    This proves the case for $k=1$.
    Next we perform the induction step.
    Assume that the lemma holds for some $k \in \{1, 2, \ldots, \ell-1\}$.
    The check of \cref{enu:hiding-same} is executed with probablity
    $\Omega(1/\ell)$; using that $\strategy'$ succeeds in \cref{enu:hiding-same}
    (conditioned on the right types having been sampled) with probablity at least $1 - O(\ell
    \delta)$, we have
    \begin{equation*}
      \sum_{y_{\le k},\, y^\perp_{k+1},\, x_{>k+1}}
      \bra{\psi} \ahk[{y_{<k},\, [L^\perp_{k+1,\, y_{\le k}}
      (\cdot ) = y^\perp_{k+1}],\, x_{>k+1}}] \otimes \bhp[y_{\le k},\,
      y^\perp_{k+1},\, x_{>k+1}] \ket{\psi} \ge 1 - O(\ell \delta)\;.
    \end{equation*}

    We now apply \cref{lem:conditional-consistency}, choosing the measurements
    $A$, $B$ and outcomes $x$, $y$, $z$ in the lemma as follows:
    \begin{gather*}
      \text{``$A$''}: \ahk[{y_{<k},\, [L^\perp_{k+1,\, y_{\le k}}(\cdot) =
          y^\perp_{k+1}],\, x_{>k+1}}]\;, \quad
      \text{``$B$''}: \bhp[y_{\le k},\, y^\perp_{k+1},\, x_{>k+1}]\;, \\
      \text{``$x$''}: y_{<k}\;, \quad
      \text{``$y$''}: y_k\;, \quad
      \text{``$z$''}: (y^\perp_{k+1},\, x_{>k+1})\;.
    \end{gather*}
		Note that here $A$ does not depend on $y$, so we use the same $A$ for all values of $y$. 
    \Cref{lem:conditional-consistency} with the above choices of parameters
    implies that
    \begin{equation*}
      \begin{split}
        \ia \otimes \bhp[y_{\le k},\, y^\perp_{k+1},\, x_{>k+1}]
        & \approx_\delta \ahk[{y_{<k},\, [L^\perp_{k+1,\, y_{\le k}}(\cdot) =
          y^\perp_{k+1}],\, x_{>k+1}}] \otimes \bhp[y_{\le k}]\\
        & \approx_\delta  \bigl( \zLlk \otimes \xLperpk[k+1] \otimes \x[x_{>k+1}]
        \bigr) \otimes \bhp[y_{\le k}]\\
        & \approx_\delta \bigl( \zLlk \zLlek \otimes \xLperpk[k+1] \otimes \x[x_{>k+1}]
        \bigr) \otimes \ib,\\
        & \approx_0 \bigl( \zLlk[k+1] \otimes \xLperpk[k+1] \otimes \x[x_{>k+1}]
        \bigr) \otimes \ib,
      \end{split}
    \end{equation*}
    where the input to $L^\perp_{k+1,\, y_{<k+1}}(\cdot)$ is $x_{k+1}$.
    The second approximation uses the induction hypothesis,
    \cref{fact:data-processing}, \cref{fact:agreement}, and
    \cref{fact:add-a-proj}.
    The third approximation follows from \cref{eq:intro-sound-4a-2} and
    \cref{fact:add-a-proj}.
    The fourth approximation follows from the definition of CL functions.
    This completes the induction.
  \end{proof}

  \begin{lemma}
    For all $k \in \{1, \ldots, \ell\}$,
	  \begin{equation}
      \label{eq:intro-sound-4}
      \ia \otimes \br[y_{<k},\, y_k^\perp] \approx_\delta
      \bigl( \zLlk \otimes \xLperpk \bigr) \otimes \ib \;.
    \end{equation}
    Moreover, analogous equations hold with operators acting on the other side
    of the tensor product.
  \end{lemma}

  \begin{proof}
    \Cref{lem:hide-rigidity,fact:data-processing} imply that
    \begin{equation}\label{eq:intro-sound-3d}
      \ahk[y_{<k},\, y^\perp_k] \otimes \ib \approx_\delta
      \Bigl( \zLlk \otimes \xLperpk \Bigr) \otimes \ib\:.
    \end{equation}
    Since the strategy $\strategy'$ succeeds in Item~\ref{enu:hiding-read} with
    probability at least $1 - O(\ell \delta)$ it follows from
    \Cref{fact:data-processing} that
    \begin{equation}\label{eq:intro-sound-3e}
      A^{\Hide{\ell}}_{y_{<\ell},\, y^\perp_{\leq \ell}} \otimes \ib
      \approx_\delta \ia \otimes \br[y_{<\ell},\, y^\perp_{\leq
          \ell}].
    \end{equation}
    An inductive argument applied to~\eqref{eq:hide-chain} of
    Lemma~\ref{lem:intro-sound-2}, combined with \Cref{fact:data-processing},
    implies that for all $1 \leq k \leq \ell$ we have
    \begin{equation}\label{eq:intro-sound-3f}
      A^{\Hide{k}}_{y_{<k},\, y^\perp_{k}} \otimes \ib \approx_\delta A^{\Hide{
          \ell}}_{y_{<k},\, y^\perp_{k}} \otimes \ib.
    \end{equation}
    Equations~\eqref{eq:intro-sound-3d},~\eqref{eq:intro-sound-3e},
    and~\eqref{eq:intro-sound-3f}, combined with \Cref{fact:data-processing},
    then establishes the lemma statement.
  \end{proof}

  \begin{lemma}
    \label{lem:intro-sound-induction}
    For each $k\in \{1, 2, \ldots, \ell+1\}$, there exists a product state
    $\ket{\ancilla_k} = \ket{\ancilla_{k,\alice}} \otimes
    \ket{\ancilla_{k,\bob}} \in \mH_{\alice'_k} \otimes \mH_{\bob'_k}$ and, for
    each $y_{<k} \in L_{<k}(V)$, a projective measurement $\bigl\{\aigek
    \bigr\}$ that acts on $\mH_{\alice} \otimes \mH_{\alice'_k}$ such that the
    following holds.
		First, for all $y \in V$, the operator $\aigek$ acts as identity on the
    register subspace spanned by basis vectors for the subspace $V_{<
      k}(y_{<k})$, and as a consequence the operator
    \begin{equation}\label{eq:intro-sound-4aa}
      \aizk = \zLlk \otimes \aigek\;
    \end{equation}
    is well-defined.
    Second, let $\strategy''_k$ be the strategy defined as follows.
    The state is $\ket{\EPR_2}^{\otimes Q} \otimes \ket{\aux} \otimes
    \ket{\ancilla_k}$.
    The measurements are identical to those in $\strategy'$ defined in
    Lemma~\ref{lem:intro-pauli-strat}, except that $\bigl\{ \ai \bigr\}$ is
    replaced with $\bigl\{ \aizk \bigr\}$.
    Then $\strategy''_k$ succeeds with probability at least $1-\delta$ in the
    game $\verifier_n^\intro$.
  \end{lemma}
  
  \begin{proof}

    The proof is by induction on $k$ from $1$ to $\ell+1$.
    The case $k=1$ is trivial by setting $A_{y_{\geq 1},\, a}^{\Intro,\, y_{<1}}
    = \ai[y,\, a]$ for all $y,a$.
    Assume that for some $k \in \{1, 2, \ldots, \ell\}$ there exist projective
    measurements $\bigl\{\aigek \bigr\}$ for every $y_{< k}$ and a strategy
    $\strategy''_k$ satisfying the conditions of the lemma statement.
    We show the statement of the lemma for $k+1$.

    \paragraph{Commutation with $Z$-basis measurements.}
    We first prove that on average over $y_{< k}$, the measurement operator
    $\aigek$, which comes from the inductive assumption, commutes with the
    projective measurement $\{\sigma^Z_{z_k}\}$ where the outcomes $z_k$ range
    over the factor space $V_k(y_{<k})$.

    To do so, we first apply Lemma~\ref{lem:commutation-analysis} where we
    choose the measurements ``$A$'', ``$B$'', and~``$C$'' and outcomes ``$a$'',
    ``$b$'', and~``$c$'' in the lemma as follows:
    \begin{gather*}
      \text{``$A$''}: \{\aizk\}\;,\quad
      \text{``$B$''}: \{\bs[z, \, a]\}\;,\quad
      \text{``$C$''}: \{\z  \}\;, \\
      \text{``$a$''}: y_{< k} \;, \quad
      \text{``$b$''}: (y_{\geq k},a) \;, \quad
      \text{``$c$''}: z \;.
    \end{gather*}
    To make sense of how the ``$B$'' POVM is indexed by ``$a$'', ``$b$'', and
    ``$c$'' as described above, we use the following relabelling: for all
    $(z,a)$, identify $\bs[z,\, a]$ with $\bs[y,\, a,\, z]$ where $y =
    L(z)$.
    Similarly, for the ``$C$'' POVM, we identify $\sigma^Z_z$ with the operator
    $\sigma^Z_{y_{<k},\, z}$ where $y_{<k} = L_{<k}(z)$.
    By applying Lemma~\ref{lem:intro-sound-1} to $\strategy_k''$ (the strategy
    given by the inductive hypothesis) with ``$k$'' in
    Lemma~\ref{lem:intro-sound-1} set to $\ell$, we have that
    \begin{equation}
      \label{eq:intro-sound-bsL2}
      \aizk \otimes \ib \approx_\delta \ia \otimes \bs[y,\, a]\;
    \end{equation}
    where $\bs[y,\, a] = \sum_{z: L(z) = y} \bs[z,\, a]$.
    Equations~\eqref{eq:intro-sound-1} and~\eqref{eq:intro-sound-bsL2} imply that the
    conditions of Lemma~\ref{lem:commutation-analysis} are satisfied, and thus
    we obtain
    \begin{equation}\label{eq:commute-with-Z}
      \bigl[\aizk \,,\, \z \bigr] \otimes \ib \approx_\delta 0\;
    \end{equation}
    where in the answer summation, $y$ is a deterministic function of $z$.
    We now apply Lemma~\ref{lem:commutation-simplification}, choosing the
    measurements ``$A$'', ``$B$'', ``$Q$'' and outcomes ``$a$'', ``$b$'',
    ``$c$'' in the lemma as follows:
    \begin{gather*}
      \text{``$A$''}: \{\aizk\}\;,\quad
      \text{``$B$''}: \{\z[\Llk] \otimes \z[z_k] \}\;, \quad
      \text{``$Q$''}: \{\z[\Llk] \}\;, \\
      \text{``$a$''}: y_{< k} \;, \quad
      \text{``$b$''}: (y_{\geq k}, a) \;, \quad
      \text{``$c$''}: z_k \;.
    \end{gather*}
    \hnote{Added in response to Bowen comments, second file, item 20} Here, we write $\z[\Llk] \otimes \z[z_k]$ to denote
    \[
    	\sum_{\substack{z' \in V: \\ L_{<k}(z') = y_{<k} \\ (z')^{V_k(y_{<k})} = z_k}} \sigma^Z_{z'}.
    \]
    We choose the Hilbert spaces ``$\mH_{\alice,a}$'' and ``$\mH_{\bob,a}$'' as
    the register subspace $ V_{<k}(y_{<k})$, and ``$\mH_{\alice',a}$'' and
    ``$\mH_{\bob',a}$'' as the register subspace $V_{\geq k}(y_{<k})$ tensored
    with $\mH_{\alice''} \otimes \mH_{\bob''}$, the Hilbert space of the state
    $\ket{\aux}$.
    Thus for every $y_{<k}$, the state $\ket{\EPR}^{\otimes Q} \otimes
    \ket{\aux}$ of the strategy $\strategy_k''$ can be decomposed into a tensor
    product of a ``question state'' and an ``answer state'' as follows:
    \[
      \Big ( \ket{\EPR_2}_{V_{<k}(y_{<k})} \Big)_{\mH_{\alice,a} \mH_{\bob,a}}
      \otimes \Big ( \ket{\EPR_2}_{V_{\geq k}(y_{<k})} \otimes \ket{\aux}
      \Big)_{\mH_{\alice',a} \mH_{\bob',a}}\;.
    \]

    Let $\mu_{L_{<k}}$ denote the distribution over outcomes $y_{< k}$ generated
    by performing the ``$Q$'' measurement on the state $\ket{\EPR}^{\otimes Q}$,
    which is equivalent to the distribution generated by the following
    procedure: (1) sample a uniformly random $z \in V$; (2) compute $y = L(z)$;
    (3) return $y_{<k}$.
    Then, since Equation~\eqref{eq:commute-with-Z} and the inductive hypothesis
    about the structure of $A^{\Intro\,,Z_{<k}}_{y_{\leq k},\,a}$ satisfy the
    conditions of Lemma~\ref{lem:commutation-simplification}, we obtain on
    average over $y_{< k} \sim \mu_{L_{< k}}$
    \begin{equation}
      \label{eq:intro-sound-5}
      \bigl[ \aigek \,,\, \z[z_k] \bigr] \otimes \ib \approx_\delta 0 \;.
    \end{equation}
    Here, the measurement outcomes $z_k$ range over $V_k(y_{<k})$.

    \paragraph{Commutation with $X$-basis measurements.}
    Next, we first prove that on average over $y_{< k}$, the measurement
    operator $\aigek$ commutes with the projective measurement
    $\bigl\{ \xLperpk \bigr\}$ where the
    outcomes $y_k^\perp$ are elements of the factor space $V_k(y_{<k})$.
  
    We again apply Lemma~\ref{lem:commutation-analysis}, choosing the
    measurements and outcomes in the lemma as follows:
    \begin{gather*}
      \text{``$A$''}: \{\aizk\}\;,\quad
      \text{``$B$''}: \{ \br[y,\, y^\perp_k,\, a] \}\;,\quad
      \text{``$C$''}: \{ \z[\Llk] \otimes \xLperpk \}\;, \\
      \text{``$a$''}: y_{< k} \;, \quad
      \text{``$b$''}: (y_{\geq k},a) \;, \quad
      \text{``$c$''}: y_k^\perp \;.
    \end{gather*}
    Equations~\eqref{eq:intro-sound-4letter-before-a} (with ``$k$'' in
    Lemma~\ref{lem:intro-sound-2} chosen to be $\ell+1$)
    and~\eqref{eq:intro-sound-4} imply the conditions of
    Lemma~\ref{lem:commutation-analysis}, so we obtain
    \begin{equation}\label{eq:commute-with-some-X}
      \Bigl[ \aizk \,,\,  \zLlk \otimes \xLperpk \Bigr] \otimes
      \ib \approx_\delta 0\;.
    \end{equation}
    We then apply Lemma~\ref{lem:commutation-simplification} with the following
    choice of measurements and outcomes:
    \begin{gather*}
      \text{``$A$''}: \{\aizk\}\;,\quad
      \text{``$B$''}: \{\zLlk \otimes \xLperpk \}\;, \quad
      \text{``$Q$''}: \{\zLlk \}\;, \\
      \text{``$a$''}: y_{< k} \;, \quad
      \text{``$b$''}: (y_{\geq k},a) \;, \quad
      \text{``$c$''}: y_k^\perp \;.
    \end{gather*}
    The Hilbert space and state decomposition are the same as in the previous
    invocation of Lemma~\ref{lem:commutation-simplification}.
    Equation~\eqref{eq:commute-with-some-X} and the inductive hypothesis satisfy
    the conditions of Lemma~\ref{lem:commutation-simplification}, and we
    similarly obtain that on average over $y_{< k} \sim \mu_{L_{<k}}$,
    \begin{equation}
      \label{eq:intro-sound-6}
      \bigl[ \aigek \,,\, \xLperpk \bigr] \otimes \ib
      \approx_\delta 0\;,
    \end{equation}

    \paragraph{Applying the Pauli twirl.}
    The last step is to apply the Pauli twirl to decompose the family of
    measurements $\{\aigek \}$ into a tensor product measurement, with the first
    part of the tensor product measuring the $k$-th linear map of $L$.

    Again applying Lemma~\ref{lem:intro-sound-1} to $\strategy_k''$, and using
    Facts~\ref{fact:data-processing} and~\ref{fact:agreement}, we obtain that
    \begin{equation*}
      \aizk[y_{\le k}] \otimes \ib \approx_\delta \zLlek \otimes \ib\;
    \end{equation*}
    which is equivalent to, by the inductive hypothesis,
    \begin{equation}\label{eq:intro-sound-8}
      \zLlk \otimes \aigek[y_k] \otimes \ib \approx_\delta \zLlk \otimes \zLk
      \otimes \ib\;.
    \end{equation}
    Applying Lemma~\ref{lem:conditional} to Equation~\eqref{eq:intro-sound-8},
    we conclude that
    \begin{equation}\label{eq:consistency-equation}
      \aigek[y_k] \otimes \ib \approx_\delta
      \zLk \otimes \ib\;,
    \end{equation}
    on average over $y_{<k} \sim \mu_{L_{<k}}$.
  
    Now we apply Lemma~\ref{lem:mixing} with the following identification:
    \begin{gather*}
      \text{``$x$''}: y_{< k} \;, \quad
      \text{``$y$''}: y_{k} \;, \quad
      \text{``$a$''}: (y_{> k}, a) \;, \quad
      \text{``$L_x$''}: L_{k,\, y_{<k}} \;,\\
      \text{``$U_x$''}: V_k(y_{< k}) \;, \quad
      \text{``$M^{x}_{y,\, a}$''}: \aigek \;, \quad
      \text{``$M^{x,\, y}_a$''}: \aigk \;.
    \end{gather*}
    The ``Consistency'' condition is implied by
    Equation~\eqref{eq:consistency-equation} and the ``Commutation'' conditions
    are implied by Equations~\eqref{eq:intro-sound-5}
    and~\eqref{eq:intro-sound-6}.
    We obtain that for all $y_{\leq k}$ there exists POVM measurements $\bigl\{
    \aigk \bigr\}$ that act as identity on the register subspace spanned by
    basis vectors for the subspace $V_{< k}(y_{<k})$ and, on average over
    $y_{<k} \sim \mu_{L_{<k}}$, we have
    \begin{equation}\label{eq:intro-sound-7}
      \aigek \otimes \ib
      \approx_\delta \Bigl( \zLk \otimes \aigk  \Bigr)
      \otimes \ib\;.
    \end{equation}
    Using the inductive assumption on the structure of $A^{\intro,\,
      Z_{<k}}_{y,\, a} $ from~\eqref{eq:intro-sound-4aa}, we get
    \begin{align}
      \aizk \otimes \ib
      & = \Bigl( \zLlk \otimes \aigek \Bigr) \otimes \ib \\
      & \approx_\delta \bigl( \zLlek \otimes \aigk \bigr) \otimes \ib\;.
        \label{eq:whatever}
    \end{align}
    where the second line follows from~\eqref{eq:intro-sound-7} and
    Lemma~\ref{lem:conditional}.
    By~\eqref{eq:whatever}, replacing the projective measurement $\bigl\{\aizk
    \bigr\}$ with the POVM
    \begin{equation*}
      \bigl\{\zLlek \otimes \aigk \bigr\}
    \end{equation*}
    in the strategy $\strategy''_k$ results in a strategy that succeeds with
    probability at least $1 - \delta$.
    To show that $\{\aigk\}$ can furthermore be turned into a projective
    measurement we use Naimark's theorem as formulated in~\cite[Theorem
    5.1]{ML20}. In particular, we see that the state
    $\ket{\ancilla'}$ it produces is a product state with no entanglement.
    This yields a strategy $\strategy_{k+1}''$ with ancilla state
    $\ket{\ancilla_{k+1}} = \ket{\ancilla_k}\ket{\ancilla'}$, which establishes
    the induction hypothesis for $k+1$.
    This completes the proof of Lemma~\ref{lem:intro-sound-induction}.
  \end{proof}

  Taking $k = \ell+1$ in Lemma~\ref{lem:intro-sound-induction} we obtain a
  strategy $\strategy''=\strategy''_{\ell+1}$ with value~$1-\delta$ in which,
  given question $(\Introspect, \alice)$, player~$\alice$ performs the
  measurement
  \begin{equation}
    \label{eq:intro-sound-final}
    \bigl\{ \z[{[L^\alice(\cdot)=y]}] \otimes A^{\Intro,\, y}_a \bigr\} \;,
  \end{equation}
  for a family of measurements $\bigl\{ A^{\Intro,\, y}_a \bigr\}$ acting on the
  state $\ket{\EPR}_{\overline{V}} \otimes \ket{\aux} \otimes \ket{\ancilla}$
where $\ket{\ancilla}=\ket{\ancilla_{\ell+1}}$ is unentangled.
  An analogous argument for Player~$\bob$'s measurements shows that we may
  additionally assume Player~$\bob$ responds to the question~$(\Introspect,
  \bob)$ using the measurement
  \begin{equation}
    \label{eq:intro-sound-final-bob}
    \bigl\{\sigma^Z_{[L^\bob(\cdot)=y]} \otimes
    B^{\Intro,\, y}_a\bigr\}\;,  
  \end{equation}
  for a family of measurements $\bigl\{ B^{\Intro,\, y}_a \bigr\}$ acting on the
  state $\ket{\EPR}_{\overline{V}} \otimes\ket{\aux} \otimes \ket{\ancilla'}$,
  where $\ket{\ancilla'}$ is unentangled.
  Summarizing, the strategy $\strategy''$ uses the state
  \begin{equation*}
    \ket{\theta} = \ket{\EPR}^{\otimes Q} \otimes \ket{\aux} \otimes
    \ket{\ancilla} \otimes \ket{\ancilla'}\;,
  \end{equation*}
  and the measurements given by Equations~\eqref{eq:intro-sound-final}
  and~\eqref{eq:intro-sound-final-bob}.
  
	To conclude the proof of the soundness part of the theorem we analyze
  Item~\ref{enu:intro-game} in Figure~\ref{fig:intro-decider}.
  The test in Item~\ref{enu:intro-game} is executed with probability
  $\Omega(1/\ell)$, so the strategy $\strategy''$ succeeds with probability at
  least $1 - \delta$ in that test, conditioned on the right types (here we
  absorb factors of $O(\ell)$ into $\delta$).
  Using~\eqref{eq:intro-sound-final} and~\eqref{eq:intro-sound-final-bob},
  conditioned on the test being executed the distribution of the
  part~$(y_\alice, y_\bob)$ of the players' answers in the test is exactly the
  distribution $\mu_{\sampler, N}$ associated with game~$\verifier_N$.
  As a result, the strategy which uses the state $\ket{\EPR}_{\overline{V}}
  \otimes \ket{\aux} \otimes \ket{\ancilla} \otimes \ket{\ancilla'}$ and
  measurements $ \{ A^{\Intro,\, y}_a \},\{ B^{\Intro,\, y}_a \}$ succeeds with
  probability at least~$1- \delta$ in the game~$\verifier_N$.
  Thus, $\val^*(\verifier_N) \geq 1- \delta$.
  It remains to show that $\delta$ has the required form. 
  As $\delta = \poly(\delta')$, $\delta'(\eps, R) = a' \bigl( (\log
  R)^{a'}\eps^{b'} + (\log R)^{-b'} \bigr)$, and $R = (2^n)^\lambda$, there is a
  constant $C \ge 1$ such that
  \begin{equation*}
    \begin{split}
      \delta(\eps, n) & \le C \Bigl( a' \bigl( (\lambda n)^{a'} \eps^{b'} +
      (\lambda n)^{-b'} \bigr) \Bigr)^{\frac{1}{C}}\\
      & \le C (a')^{1/C} \bigl(  (\lambda n)^{a'/C} \cdot \eps^{b'/C} + (\lambda n)^{-b'/C} \bigr)\\
      & \le a \bigl( (\lambda n)^a \cdot \eps^b + (\lambda n)^{-b} \bigr)\;,
    \end{split}
  \end{equation*}
  for $a = \max \bigl\{ C (a')^{1/C}, a'/C \bigr\}$, and $b = b'/C$, where the
  second line uses the inequality $(x+y)^{1/C} \le x^{1/C} + y^{1/C}$ for
  $x,y>0$ and $C\ge 1$. 
  
  
  This establishes the soundness part of the theorem.

  To show the entanglement bound, we observe that local isometries do not change
  the Schmidt rank of a state.
  Define $\ket{\psi'} = \phi(\ket{\psi}) \otimes \ket{\ancilla} \otimes
  \ket{\ancilla'}$.
  Since the strategy $(\ket{\psi'}, \{ A^{\Intro,\, y}_a \},\{ B^{\Intro, y}_a
  \})$ is $\delta$-close to $(\ket{\theta}, \{ A^{\Intro,\, y}_a \},\{
  B^{\Intro,\, y}_a \})$ which has value $1 - \delta$, the strategy
  $(\ket{\psi'}, \{ A^{\Intro,\, y}_a \},\{ B^{\Intro,\, y}_a \})$ has value at
  least $1 - 2\delta$ in the game $\verifier_N$, and therefore the Schmidt rank
  of $\phi(\ket{\psi})$ (and thus of $\ket{\psi}$) must be at least
  $\Ent(\verifier_N,1 - 2\delta)$.
  Here, we use the fact that ancilla states are product states and therefore
  have Schmidt rank~$1$.
  
  Moreover, recall that $\phi(\ket{\psi})$ is $\delta$-close to
  $\ket{\EPR_2}^{\otimes Q} \otimes \ket{\aux}$ whose Schmidt coefficients are
  all at most $2^{-Q/2}$.
  For any bipartite state $\ket{a}$ with Schmidt rank at most $R$ and $\ket{b}$
  whose Schmidt coefficients are all at most $\beta$, it follows from the
  Cauchy-Schwarz inequality that $|\langle a | b \rangle|^2 \leq R \beta^2$.
  Therefore the Schmidt rank of $\phi(\ket{\psi})$ (and thus of $\ket{\psi}$) is
  at least
  \[
  	\left (1 - \delta \right)^2 \cdot 2^{Q} \geq \left
      (1 - \delta \right)^2 \cdot 2^{N^\lambda}\;,
  \]
  where we used that $Q = M \log q \geq R$ (using the canonical parameter
  settings of Definition~\ref{def:introparams}) and $\delta \geq \|
  \phi(\ket{\psi}) - \ket{\theta} \|^2 = 2 - 2\Re \bra{\theta}
  \phi(\ket{\psi})$.
  Combining the two lower bounds on the Schmidt rank of $\ket{\psi}$ shows the
  desired lower bound on $\Ent(\tvint_n,1 - \eps)$.
\end{proof}

 
\section{Oracularization}
